% Supplement section S5: elder selection.
\section{Elder selection}\label{sec:selection}

This section concerns the selection bound (3.13) of the main text: under the determinant-weighted Palm law of a typed maximum--index-two-saddle contact pair at separation $r$, the probability that the pair is not the elder persistence pair is at most $C_{\rm sel}r^3$, uniformly on compact positive marks. The main text takes (3.13) from Theorem~A of \cite{parentP} (Theorem~\ref{thm:parent-A} of this supplement), which proves it for every $d\ge2$ and $L>0$ (Remark~\ref{rem:selection-parent}). This section proves the statements that the main text quotes for the Palm law and the elder mark (Definitions~\ref{def:selection-palm}, \ref{def:selection-mark} and~\ref{def:selection-prob}, Lemma~\ref{lem:selection-mark}) and the trichotomy (Lemma~\ref{lem:selection-trichotomy}). It then gives a second, three-dimensional proof of (3.13) through explicit witnesses (Theorem~\ref{thm:selection-main}): the regional witness counts behind (3.14)--(3.17) of the main text (\S\ref{ssec:selection-witness}) and the capture--escape theorem (Theorem~\ref{thm:selection-escape}).

Throughout, $f=f_{24}$ is the normalized side-$24$ periodized Bargmann--Fock field on $\T^3_{24}$, $0<r<12$, $t\in S^2$, $b\in\R$ and $\kappa>0$; the pair is $M=x-\frac r2t$, $S=x+\frac r2t$, with midpoint $x$, and $h=f(S)$. We use an orthonormal frame $(e_1,e_2,e_3)=(t,v,w)$ whose first vector is $t$, and write $f_{ij\cdots}$ for partial derivatives in this frame. Constants $C,c$ depend only on the compact mark set and on the fixed charts, and may change from line to line.

\subsection{The Palm law and the elder mark}\label{ssec:selection-palm}

\begin{definition}[pins, weight and Palm law]\label{def:selection-palm}
For $0<r<12$, $t\in S^2$, $b\in\R$ and $\kappa>0$ let
\[
\Pi_{r,x,t}=\bigl(f(M),f(S),\nabla f(M),\nabla f(S)\bigr),\qquad z_r(b,\kappa)=\bigl(b,\,b-\kappa r^3/6,\,0,\,0\bigr).
\]
The conditional law $\E^0_{r,t,b,\kappa}[\,\cdot\,]=\E[\,\cdot\mid\Pi_{r,x,t}=z_r(b,\kappa)]$ is the regular conditional Gaussian distribution. It exists because the pin vector is Gaussian, and since the corrected-pin covariance is nondegenerate (Section~\ref{sec:setting}), the conditional distributions have a version that is continuous in $(r,t,b,\kappa)$. Put
\[
W_r=|\det H_M|\,|\det H_S|\,\ind_{\{H_M\prec0,\ \Ind(H_S)=2\}},\qquad Z_r(t,b,\kappa)=\E^0_{r,t,b,\kappa}[W_r].
\]
By Section~\ref{sec:contact}, $c_Zr^2\le Z_r\le C_Zr^2$ on compact positive mark sets; in particular $Z_r>0$. The \emph{all-typed determinant-weighted Palm law} is
\begin{equation}\label{eq:selection-palm}
\E^{MS}_{r,t,b,\kappa}[F]=\frac{\E^0_{r,t,b,\kappa}[W_rF]}{Z_r(t,b,\kappa)} .
\end{equation}
We write $\Prob^{MS}_r$ and $\Prob^0$ for the corresponding probabilities, and drop the subscripts $t,b,\kappa$ when they are fixed.
\end{definition}

\begin{definition}[maximin level and elder mark]\label{def:selection-mark}
Let $\Omega_{\rm Mo}$ be the event that $f$ is Morse with pairwise distinct critical values; $\Prob(\Omega_{\rm Mo})=1$ by Proposition~1.1 of the main text. For continuous $g:\T^3_{24}\to\R$ and $x\in\T^3_{24}$, the \emph{maximin connection level} of $x$ is
\begin{equation}\label{eq:selection-connection}
d_g(x)=\sup\Bigl\{\min_{s\in[0,1]}g(\gamma(s))\ :\ \gamma:[0,1]\to\T^3_{24}\ \text{continuous},\ \gamma(0)=x,\ g(\gamma(1))>g(x)\Bigr\},
\end{equation}
with $\sup\emptyset=-\infty$. Since $\T^3_{24}$ is path-connected, $d_g(x)=-\infty$ exactly when $x$ is a global maximum of $g$. For a typed pair $(M,S)$ with $h=f(S)<f(M)$, the \emph{elder mark} is
\begin{equation}\label{eq:selection-mark}
e(f;M,S)=\ind\{d_f(M)=f(S)\}.
\end{equation}
\end{definition}

\begin{lemma}[the mark is correct and Borel]\label{lem:selection-mark}
On $\Omega_{\rm Mo}$, let $M$ be a local maximum and $S$ an index-two critical point with $h=f(S)<f(M)$. Then $e(f;M,S)=1$ if and only if
\begin{enumerate}[(i)]
\item the two upper local sectors of $S$ lie in distinct components of $\{f>h\}$;
\item $M$ is the birth maximum of one of them, that is, the maximum of $f$ over that component; and
\item the birth value at $M$ is the smaller of the two.
\end{enumerate}
Conditions (i)--(iii) are exactly the elder-pair condition $S=D_f(M)$ of Section~1 of the main text. Moreover, $(g,x)\mapsto d_g(x)$ is lower semicontinuous on $C(\T^3_{24})\times\T^3_{24}$, so $(f,M,S)\mapsto e(f;M,S)$ is Borel.
\end{lemma}

\begin{proof}
Let $\mathcal U_1$ be the component of the open set $\{f>h\}$ that contains $M$. A path from $M$ that stays in $\{f>h\}$ stays in $\mathcal U_1$, and a compact path in $\{f>h\}$ has minimum $>h$.

\emph{(i)--(iii) imply $e=1$.} Let $\mathcal U_2\ne\mathcal U_1$ be the component containing the other upper sector of $S$, with birth value $b_2>f(M)$. Going from $M$ inside $\mathcal U_1$ to the first sector of $S$, through $S$, and inside $\mathcal U_2$ to its birth maximum gives a path to a point higher than $f(M)$ whose minimum is $f(S)=h$; hence $d_f(M)\ge h$. If some path from $M$ to a point higher than $f(M)$ had minimum $>h$, it would lie in $\mathcal U_1$, contradicting (ii). Hence $d_f(M)=h=f(S)$.

\emph{$e=1$ implies (i)--(iii).} Suppose $d_f(M)=h$. A point of $\mathcal U_1$ higher than $f(M)$ would give a path in $\mathcal U_1$ with minimum $>h$, so $f\le f(M)$ on $\mathcal U_1$ and $M$ is the birth maximum of $\mathcal U_1$. For $t<h$ close to $h$, the definition of $d_f(M)$ gives a path from $M$ to a point higher than $f(M)$ with minimum $>t$, so the component of $M$ in $\{f>t\}$ contains a point higher than $f(M)$, while for $t\in(h,f(M))$ it does not. Since $h$ is the only critical value in a neighbourhood of $h$ and $S$ is the only critical point at level $h$, the components of $\{f>t\}$ change between $t>h$ and $t<h$ only by attaching the one-handle at $S$, which joins the components of its two upper sectors. Hence the two sectors lie in distinct components, one of which is $\mathcal U_1$ and the other, $\mathcal U_2$, contains a point higher than $f(M)$; this is (i)--(iii).

\emph{Measurability.} For a fixed continuous path $\gamma$ from $0$, the functional $g\mapsto\min_s g(\gamma(s))$ on the open set $\{g(\gamma(1))>g(0)\}$, extended by $-\infty$ elsewhere, is lower semicontinuous in the uniform norm, and $d_g(0)$ is the supremum of these functionals over $\gamma$; a supremum of lower semicontinuous functions is lower semicontinuous. Since $d_g(x)=d_{g(\cdot+x)}(0)$ and $(g,x)\mapsto g(\cdot+x)$ is continuous, $(g,x)\mapsto d_g(x)$ is lower semicontinuous, hence Borel, and so is $\{d_f(M)=f(S)\}$.
\end{proof}

\begin{definition}[selection probability]\label{def:selection-prob}
Let $p_r(t,b,\kappa)=\E^{MS}_{r,t,b,\kappa}[e(f;M,S)]$. With $p_\Pi$ the Gaussian density of $\Pi_{r,x,t}$ and $K^{\rm all}_r=p_\Pi(z_r(b,\kappa))\,\E^0[W_r]$, the selected contact intensity is
\begin{equation}\label{eq:selection-Ksel}
K^{\rm sel}_r(t,b,\kappa)=p_r(t,b,\kappa)\,K^{\rm all}_r(t,b,\kappa).
\end{equation}
Neither adjacency nor selection enters the base law \eqref{eq:selection-palm}.
\end{definition}

\subsection{The trichotomy}\label{ssec:selection-trichotomy}

\begin{definition}[selection events]\label{def:selection-events}
On $\Omega_{\rm Mo}$ fix an ordered typed pair $(M,S)$ at distance $r$, and let $h=f(S)$ and $b=f(M)$.
\begin{enumerate}[(a)]
\item $A_r$ (\emph{adjacency}): at least one unstable half-branch of $S$ under gradient ascent converges to $M$. This is a Boolean mark: if both half-branches converge to $M$, the pair still counts once (Section~\ref{sec:setting}).
\item $E_r=\{e(f;M,S)=1\}$ (\emph{actual elder selection}, \eqref{eq:selection-mark}).
\item $L_{{\rm same},r}$ (\emph{same-maximum self-attachment}): both unstable half-branches of $S$ under gradient ascent converge to $M$.
\item $B_r$ (\emph{third-critical witness}): there is a critical point $y\notin\{M,S\}$ of $f$ with $f(y)\in[h,b]$ (value in the pair window). Its distance from the pair midpoint $x$ places it in exactly one of the regions
\begin{enumerate}[(i)]
\item \emph{collar}: $\dist(y,x)\le C_1r$;
\item \emph{singular-near}: $C_1r<\dist(y,x)\le\delta$ (so that $\eta=r/\dist(y,x)<1/C_1<1$, a compact $\eta$-chart);
\item \emph{fixed-distance}: $\dist(y,x)>\delta$, and then $\dist(y,\{M,S\})>\delta-r/2\ge\delta/2$ for $r\le\delta$.
\end{enumerate}
Here $C_1>1$ and $\delta>0$ are fixed constants of the chart decomposition, and $r\le\delta/C_1$.
\end{enumerate}
\end{definition}

\begin{lemma}[merge level]\label{lem:selection-merge}
On $\Omega_{\rm Mo}$, let $M\ne M'$ be local maxima of $f$ that lie in one component of $\{f>h\}$, and let
\[
\tau=\sup\Bigl\{\min_{s\in[0,1]}f(\gamma(s))\ :\ \gamma:[0,1]\to\T^3_{24}\ \text{continuous},\ \gamma(0)=M,\ \gamma(1)=M'\Bigr\}.
\]
Then $h<\tau<\min(f(M),f(M'))$, and $\tau=f(y)$ for a critical point $y$ of $f$. In particular $y\notin\{M,M',S\}$, and $f(y)\in(h,f(M))$.
\end{lemma}

\begin{proof}
A path from $M$ to $M'$ inside their common component of $\{f>h\}$ has compact image, hence minimum larger than $h$; so $\tau>h$. Since $M$ is a strict local maximum, there is a small sphere around $M$ on which $f\le f(M)-\varepsilon$ for some $\varepsilon>0$, and every path from $M$ to $M'$ crosses it; so $\tau\le f(M)-\varepsilon$, and likewise for $M'$. Suppose that $\tau$ is not a critical value. The critical values are finite in number, so there is $\varepsilon>0$ with no critical value in $[\tau-\varepsilon,\tau+\varepsilon]$ and with $\tau+\varepsilon<\min(f(M),f(M'))$. The flow of $\nabla f/|\nabla f|^2$, stopped on reaching level $\tau+\varepsilon$, deforms $\{f\ge\tau-\varepsilon\}$ into $\{f\ge\tau+\varepsilon\}$ and fixes $\{f\ge\tau+\varepsilon\}$, which contains $M$ and $M'$. It carries a path from $M$ to $M'$ with minimum larger than $\tau-\varepsilon$ to one with minimum at least $\tau+\varepsilon>\tau$, contradicting the definition of $\tau$. Hence $\tau=f(y)$ for a critical point $y$. Its value lies strictly between $f(S)=h$ and $\min(f(M),f(M'))$, so $y\notin\{M,M',S\}$.
\end{proof}

\begin{lemma}[exhaustive trichotomy]\label{lem:selection-trichotomy}
On $\Omega_{\rm Mo}$,
\begin{equation}\label{eq:selection-trichotomy}
A_r\cap E_r^c\subset B_r\cup L_{{\rm same},r}.
\end{equation}
\end{lemma}

\begin{proof}
Work on $A_r$. A half-branch of $S$ ascends to $M$; it leaves $S$ through one upper local sector of $S$ and stays in $\{f>h\}$, so $M$ lies in the component $\mathcal U_1$ of $\{f>h\}$ that contains this sector. Let $\mathcal U_2$ be the component containing the other upper sector. It suffices to find a critical point $y\notin\{M,S\}$ with $f(y)\in(h,b)$; the regions (i)--(iii) of Definition~\ref{def:selection-events} then place it in $B_r$.

\emph{Case 0: $\mathcal U_1=\mathcal U_2$.} The other half-branch also stays in $\mathcal U_1$ and converges to a critical point $z$ of $f$ with $f(z)>h$. If $z=M$, both half-branches converge to $M$, which is $L_{{\rm same},r}$. Otherwise $z\ne M$, and $z\ne S$ because $f(z)>h$. If $f(z)<f(M)=b$, then $y=z$ is the required point. If $f(z)>f(M)$, then $\mathcal U_1$ contains a point higher than $M$. On the boundary of $\mathcal U_1$, $f=h<f(M)$, so the maximum of $f$ over the closure of $\mathcal U_1$ is attained inside $\mathcal U_1$, at a local maximum $M'$ with $f(M')>f(M)$. Lemma~\ref{lem:selection-merge}, applied to $M$ and $M'$, gives the required point $y$.

For Cases~1 and~2, assume $\mathcal U_1\ne\mathcal U_2$, with birth values $b_1,b_2$ and birth maxima $M_1,M_2$; distinct critical values make these unique.

\emph{Case 1: $M\ne M_1$.} Then $M_1$ and $M$ are distinct maxima in the component $\mathcal U_1$ of $\{f>h\}$, and Lemma~\ref{lem:selection-merge} gives a critical point $y\notin\{M,S\}$ with $f(y)\in(h,f(M))=(h,b)$, a window witness.

\emph{Case 2: $M=M_1$ but $e=0$.} Since $\mathcal U_1\ne\mathcal U_2$ and condition~(ii) of Lemma~\ref{lem:selection-mark} holds, failure of \eqref{eq:selection-mark} forces $f(M)>\min(b_1,b_2)$, that is, $b_2<f(M)$: $\mathcal U_2$ is younger and is the component killed at $S$. The required witness is the birth maximum of $\mathcal U_2$. Indeed, $y=M_2$ is a critical point (a local maximum), and its value satisfies $h<b_2=f(M_2)<f(M)=b$; the lower bound holds because $\mathcal U_2$ is a component of the \emph{open} superlevel set $\{f>h\}$ containing an upper sector of $S$, so its birth value exceeds $h$. Moreover $M_2\notin\{M,S\}$: $M_2\ne M$ because they lie in distinct components, and $M_2\ne S$ because $f(S)=h<b_2$. So $y=M_2$ is a window witness. (The merge saddle of the surviving class of $M$ would not serve here: its value can lie below $h$, outside the window.)

Outside $L_{{\rm same},r}$, each case gives a critical point $y\notin\{M,S\}$ with value in the pair window, and the regions (i)--(iii) of Definition~\ref{def:selection-events} place it in $B_r$.
\end{proof}

\begin{corollary}[reduction of the selection loss]\label{cor:selection-reduction}
Let $N_{\rm col}$ be the number of critical points $y\notin\{M,S\}$ of $f$ in the collar region~(i), and let $N_{\rm sing}$ and $N_{\rm far}$ be the numbers of critical points $y$ with $f(y)\in[h,b]$ in the singular-near region~(ii) and in the fixed-distance region~(iii). Under the Palm law \eqref{eq:selection-palm},
\begin{equation}\label{eq:selection-reduction}
1-p_r\le\Prob^{MS}_r(A_r^c)+\Prob^{MS}_r(B_r)+\Prob^{MS}_r(L_{{\rm same},r}),
\end{equation}
and
\begin{equation}\label{eq:selection-markov}
\Prob^{MS}_r(B_r)\le\E^{MS}_rN_{\rm col}+\E^{MS}_rN_{\rm sing}+\E^{MS}_rN_{\rm far}.
\end{equation}
\end{corollary}

\begin{proof}
Since $e\in\{0,1\}$, $1-p_r=\Prob^{MS}_r(E_r^c)\le\Prob^{MS}_r(A_r^c)+\Prob^{MS}_r(E_r^c\cap A_r)$, and Lemma~\ref{lem:selection-trichotomy} gives \eqref{eq:selection-reduction}. On $B_r$ at least one of the three counts is at least one, so Markov's inequality gives \eqref{eq:selection-markov}.
\end{proof}

The three terms on the right of \eqref{eq:selection-reduction} are $O(r^3)$ uniformly on compact positive mark sets: the three regional expectations in \eqref{eq:selection-markov} by Lemmas~\ref{lem:selection-collar}, \ref{lem:selection-sing} and~\ref{lem:selection-far}, and the adjacency and self-attachment failures by Proposition~\ref{prop:selection-adjacency}, which rests on the capture--escape Theorem~\ref{thm:selection-escape}. The assembly is Theorem~\ref{thm:selection-main}.

\subsection{The endpoint factor and the Palm normalization}\label{ssec:selection-endpoint}

Recall from Section~\ref{sec:contact} that, on compact positive mark sets,
\begin{equation}\label{eq:selection-normalizer}
c_Zr^2\le Z_r=\E^0[W_r]\le C_Zr^2,\qquad\E^0[W_r^2]\le C_Wr^4 .
\end{equation}
Each Palm expectation below is an unnormalized expectation under $\E^0$ divided explicitly by $Z_r$; the lower bound in \eqref{eq:selection-normalizer} is the Palm normalization used in every witness count.

\begin{lemma}[exact endpoint factor $r^2$]\label{lem:selection-endpoint}
In the frame $(t,v,w)$ write
\begin{equation}\label{eq:selection-frame}
H_p=\begin{pmatrix}a_p&m_p^{\top}\\ m_p&Q_p\end{pmatrix},\qquad p\in\{M,S\}.
\end{equation}
Under the corrected pair pins $\nabla f(M)=\nabla f(S)=0$, $f(M)-f(S)=\kappa r^3/6$,
\begin{equation}\label{eq:selection-scaled}
a_M=r\alpha_M,\qquad m_M=r\beta_M,\qquad a_S=r\alpha_S,\qquad m_S=r\beta_S,
\end{equation}
where the scaled variables are linear functionals of the third-derivative process on the segment $\{x+\vartheta t:|\vartheta|\le r/2\}$, with conditional Gaussian moments uniformly bounded on compact marks. Moreover
\begin{equation}\label{eq:selection-softdet}
\det H_M=rD_M,\qquad\det H_S=rD_S,
\end{equation}
where $D_M,D_S$ are polynomials of degree three in the scaled endpoint variables and the transverse Hessian entries. In particular $W_r\le r^2|D_M|\,|D_S|$.
\end{lemma}

\begin{proof}
For each coordinate direction $e_i$ put $g_i(\vartheta)=\partial_if(x+\vartheta t)$. Since $g_i(-r/2)=g_i(r/2)=0$, Taylor's integral formula gives the exact identities
\begin{equation}\label{eq:selection-taylor}
\begin{aligned}
\partial_{ti}f(M)&=-\frac1r\int_{-r/2}^{r/2}\Bigl(\frac r2-\vartheta\Bigr)\,\partial_{tti}f(x+\vartheta t)\,d\vartheta,\\
\partial_{ti}f(S)&=\frac1r\int_{-r/2}^{r/2}\Bigl(\vartheta+\frac r2\Bigr)\,\partial_{tti}f(x+\vartheta t)\,d\vartheta .
\end{aligned}
\end{equation}
Thus the first row and column of each endpoint Hessian carry one exact factor $r$, which is \eqref{eq:selection-scaled}. The conditional Gaussian moments of the scaled variables are uniformly bounded on compact marks because the corrected pin covariance has a uniform eigenfloor and the side-$24$ field has finite derivative moments. Expanding \eqref{eq:selection-frame} along the first row gives \eqref{eq:selection-softdet}. Type indicators can only decrease absolute values, so $W_r\le|\det H_M|\,|\det H_S|=r^2|D_M|\,|D_S|$.
\end{proof}

This is the same endpoint soft-mode factor that yields the two-sided bound in \eqref{eq:selection-normalizer}.

\begin{lemma}[conditional polynomial moments]\label{lem:selection-moments}
Let $H$ be a Gaussian vector whose conditional law given $\mathcal A=a$ has mean $\mu(a)$ with $\|\mu(a)\|\le C(1+\|a\|)$ and covariance of operator norm at most $C$. Then for every polynomial $\Psi$ of degree at most nine,
\begin{equation}\label{eq:selection-moment9}
\E\bigl[|\Psi(H)|\bigm|\mathcal A=a\bigr]\le C_\Psi\bigl(1+\|a\|^9\bigr).
\end{equation}
The product $\det H_M\det H_S\det H_y$ of three $3\times3$ Hessian determinants is a polynomial of total degree nine in the concatenated Hessian coordinates, and the product $\det H_M\det H_S$ of two such determinants is a polynomial of degree six.
\end{lemma}

\begin{proof}
Write $H=\mu(a)+G$, where $G$ is a centred Gaussian vector, possibly degenerate, with covariance of operator norm at most $C$. Expand $\Psi$ and use uniform Gaussian moments.
\end{proof}

\subsection{Regional witness counts}\label{ssec:selection-witness}

The collar and singular-near charts are first computed at the planar Bargmann--Fock jet and then transferred to the side-$24$ field through exact factorizations of the conditional determinants. These factorizations hold for every kernel satisfying the standing hypotheses below, and their remaining coefficients are positive by the derivative nondegeneracy of Section~\ref{sec:setting}. The fixed-distance count is proved directly for the side-$24$ field.

\subsubsection{Standing hypotheses and the contact-limit pair law}

Let $K_0(z)=e^{-|z|^2/2}$ be the planar Bargmann--Fock covariance and $J_0$ its even covariance jet through order six. For a stationary even kernel $\mathcal K$ let $J_6(\mathcal K)$ be its even jet through order six, with the max norm on coordinate entries.

\begin{hypothesis}[jet ball and standing hypotheses]\label{hyp:selection-jet}
\begin{enumerate}[(1)]
\item (\emph{Jet ball.})
\begin{equation}\label{eq:selection-jetball}
B=\bigl\{J:\ \|J-J_0\|_{\max}\le10^{-110}\bigr\}.
\end{equation}
\item (\emph{Standing hypotheses.}) $\mathcal K$ is a stationary, even, real-analytic covariance kernel with $J_6(\mathcal K)\in B$; the associated centred Gaussian field has the finite-family derivative nondegeneracy of Section~\ref{sec:setting} (full positive Fourier support); and all compact charts, mark sets and constants are those of the planar charts below.
\end{enumerate}
\end{hypothesis}

\begin{remark}\label{rem:selection-side24}
The normalized side-$24$ kernel $K_{24}$ satisfies Hypothesis~\ref{hyp:selection-jet}. It has full positive Fourier support (Section~\ref{sec:setting}), so Theorem~\ref{thm:setting-gram} gives the derivative nondegeneracy, and $\|J_6(K_{24})-J_0\|_{\max}<10^{-113}$ by Lemma~\ref{lem:correction-deep}. The planar kernel $K_0$ has an everywhere positive spectral density, so a family of derivative evaluations at one point is nondegenerate for $K_0$ as soon as its polynomial symbols are linearly independent. The structural coefficients $A_{\rm col}$ and $A_{\rm sing}$ of Lemmas~\ref{lem:selection-collarfactor} and~\ref{lem:selection-singfactor} are bounded below by exact factorizations and this nondegeneracy; no derivative bound for them on the jet ball is used. The jet ball enters through Lemma~\ref{lem:selection-pinstab}.
\end{remark}

\begin{lemma}[contact-limit pair pin law at the planar jet]\label{lem:selection-pairlaw}
In the contact limit the raw pin vector of the pair may be written
\begin{equation}\label{eq:selection-P}
P=(f,f_{111},f_1,f_{11},f_2,f_{12},f_3,f_{13}),
\end{equation}
all derivatives taken at $x$, with target
\begin{equation}\label{eq:selection-Ptarget}
(b,2\kappa,0,0,0,0,0,0),
\end{equation}
up to the invertible rescaling $f_{111}\mapsto-f_{111}/12$. For the planar Bargmann--Fock covariance
\begin{equation}\label{eq:selection-detP}
\det\Cov(P)=12,
\end{equation}
so the limiting pair pin law is nondegenerate. The corrected finite-$r$ pair frame converges analytically to \eqref{eq:selection-P}, so compact covariance eigenfloors and Gaussian density envelopes follow for sufficiently small $r$. Conditioned on \eqref{eq:selection-Ptarget}, the transverse Hessian
\begin{equation}\label{eq:selection-Qlaw}
Q=\begin{pmatrix}f_{22}&f_{23}\\f_{23}&f_{33}\end{pmatrix}\qquad\text{has law}\qquad Q=-bI_2+G_2,
\end{equation}
where $G_2$ is standard GOE$_2$: its diagonal entries are $N(0,2)$, its off-diagonal entry is $N(0,1)$, and these three variables are independent.
\end{lemma}

\begin{proof}
All covariances follow from $\Cov(\partial^\alpha f,\partial^\beta f)=(-1)^{|\beta|}\partial^{\alpha+\beta}K_0(0)$, and \eqref{eq:selection-detP} and \eqref{eq:selection-Qlaw} are exact Gaussian computations from these covariances.
\end{proof}

\begin{lemma}[pair-pin stability on the jet ball]\label{lem:selection-pinstab}
Let $P_r$ be the corrected pair pin vector. At contact its planar covariance has determinant $12$, and every entry is a contraction of $J_6$ with $\ell^1$ coefficient at most $3^6=729$. A crude row-sum bound of $120$ gives
\begin{equation}\label{eq:selection-pinfloor}
\lambda_{\min}\Cov_0(P)\ge\frac{12}{120^7}>3.34\times10^{-14}.
\end{equation}
On the jet ball \eqref{eq:selection-jetball} the operator perturbation of the $8\times8$ pin covariance is below
\begin{equation}\label{eq:selection-pinpert}
8\cdot729\cdot10^{-110}<5.84\times10^{-107}.
\end{equation}
Hence the pair pin covariance, its inverse, all conditional means and all conditional covariance blocks are analytic and uniformly bounded throughout the neighbourhood. This also gives uniform Gaussian polynomial moments and density envelopes.
\end{lemma}

\begin{proof}
For a positive definite $8\times8$ matrix, $\lambda_{\min}\ge\det/\lambda_{\max}^7$, and the row-sum bound gives $\lambda_{\max}\le120$; with \eqref{eq:selection-detP} this is \eqref{eq:selection-pinfloor}. Each entry moves by at most $729\cdot10^{-110}$ on the jet ball, which gives \eqref{eq:selection-pinpert}. Since \eqref{eq:selection-pinpert} is far below \eqref{eq:selection-pinfloor}, Weyl's inequality keeps the covariance uniformly positive definite, and the remaining assertions follow from the analytic dependence of Gaussian regression on the covariance.
\end{proof}

\subsubsection{The two-dimensional triple-contact elimination}

\begin{lemma}[exact triple-contact elimination]\label{lem:selection-elimination}
Let $c,\rho\ne0$ with $c^2+\rho^2=1$, and let $\eta>0$, $\kappa>0$ and $q_M,q_S\in\R$. There is a unique real cubic polynomial $p(X,Y)$ on $\R^2$ such that $M=(-\frac12,0)$, $S=(\frac12,0)$ and $y=(c/\eta,\rho/\eta)$ are critical points, $p(M)=0$, $p(S)=-\kappa/6$, $p_{YY}(M)=q_M$ and $p_{YY}(S)=q_S$. Its Hessian determinants at the three critical points are
\begin{gather}
-\det H_M=\frac{\kappa^2(4c^2-\eta^2)^2+4\kappa\rho^2\bigl[(12c^2+\eta^2)q_M+(4c^2-\eta^2)q_S\bigr]+4\rho^4(q_M-q_S)^2}{64\,c^2\rho^2},\label{eq:selection-B1}\\
-\det H_S=\frac{\kappa^2(4c^2-\eta^2)^2-4\kappa\rho^2\bigl[(4c^2-\eta^2)q_M+(12c^2+\eta^2)q_S\bigr]+4\rho^4(q_M-q_S)^2}{64\,c^2\rho^2},\label{eq:selection-B1S}\\
\det H_y=\frac{N_y}{64\,c^2\eta\rho^2},\label{eq:selection-B1y}
\end{gather}
where
\begin{multline}\label{eq:selection-Ny}
N_y=\eta\kappa^2(4c^2-\eta^2)^2+4\kappa\rho^2(4c^2+\eta^2)\bigl[\eta(q_M-q_S)-2c(q_M+q_S)\bigr]\\
-16c\rho^4(q_M^2-q_S^2)+4\eta\rho^4(q_M-q_S)^2 .
\end{multline}
In particular, the endpoint determinants are exactly quadratic in $(\kappa,q_M,q_S)$, with the common $\kappa^2$ coefficient $(4c^2-\eta^2)^2/(64c^2\rho^2)$ and the common $\kappa$-free term $4\rho^4(q_M-q_S)^2/(64c^2\rho^2)$; the product $\det H_M\det H_S\det H_y$ has exact degree six in $\kappa$; the denominators vanish only on the strata $c=0$ and $\rho=0$ (and $\eta=0$ for $\det H_y$); and on the collision strata $2c=\pm\eta$ the $\kappa^2$ coefficient $(4c^2-\eta^2)^2$ vanishes.
\end{lemma}

\begin{proof}
The ten conditions (six station equations, two value equations and two curvature prescriptions) are linear in the ten coefficients of $p$. The determinant of this linear system is $-24c\rho^3/\eta^4\ne0$, so the solution is unique. Substituting it into the Hessians at $M$, $S$ and $y$ gives \eqref{eq:selection-B1}--\eqref{eq:selection-Ny} by direct exact computation, and the stated structure is read off from these formulas.
\end{proof}

In the witness charts below, the plane is spanned by the pair axis and the direction of the third point, the chart is centred at the pair midpoint, lengths are measured in units of $r$ and values in units of $r^3$ (so that the pair sits at $(\mp\frac12,0)$ and the value gap is $\kappa/6$), $\eta=r/\dist(y,x)$, and $q_M,q_S$ are the in-plane transverse curvatures at the endpoints; in the application they are the $O(1)$ contact-law blocks. Physical two-dimensional Hessian determinants are therefore $r^2$ times the chart determinants \eqref{eq:selection-B1}--\eqref{eq:selection-B1y}. In the chart the endpoint determinants are generically of order $\kappa^2$ as $\eta\to0$. The collar count below uses these expressions only where $|c|$ is bounded below. The singular-near envelope (Lemma~\ref{lem:selection-envelope}) is proved directly from Taylor's formula and does not use them.

\subsubsection{The collar count}

\begin{lemma}[exact collar scaled-gradient chart at the planar jet]\label{lem:selection-collarchart}
Let a third point be
\begin{equation}\label{eq:selection-collarpt}
y=x+r(Xe_1+Ye_2+Ze_3),\qquad\rho^2=Y^2+Z^2,
\end{equation}
and define the physical scaled third gradient
\begin{equation}\label{eq:selection-Gcol}
G_{\rm col}=\bigl(r^{-2}f_1(y),\,r^{-1}f_2(y),\,r^{-1}f_3(y)\bigr).
\end{equation}
The pair pins make \eqref{eq:selection-Gcol} converge to $(L_1,L_\perp)$, where
\begin{gather}
L_1=\tfrac12\bigl(X^2-\tfrac14\bigr)f_{111}+XYf_{112}+XZf_{113}+\tfrac{Y^2}2f_{122}+YZf_{123}+\tfrac{Z^2}2f_{133},\label{eq:selection-L1}\\
L_\perp=Q(Y,Z)^{\top}.\label{eq:selection-Lperp}
\end{gather}
The axial and transverse blocks are conditionally independent. Their exact conditional means are
\begin{equation}\label{eq:selection-colmeans}
\E[L_1\mid P]=\kappa\bigl(X^2-\tfrac14\bigr),\qquad\E[L_\perp\mid P]=-b(Y,Z)^{\top},
\end{equation}
the exact conditional axial variance is
\begin{equation}\label{eq:selection-colvar}
\Var(L_1\mid P)=\frac{\rho^2(4X^2+\rho^2)}2,
\end{equation}
and the transverse covariance is
\begin{equation}\label{eq:selection-Cperp}
C_\perp=\begin{pmatrix}2Y^2+Z^2&YZ\\YZ&Y^2+2Z^2\end{pmatrix},\qquad\det C_\perp=2\rho^4,\qquad(Y,Z)\,C_\perp^{-1}(Y,Z)^{\top}=\tfrac12 .
\end{equation}
Therefore
\begin{equation}\label{eq:selection-coldet}
\det\Cov\bigl(\lim G_{\rm col}\bigm|P\bigr)=\rho^6(4X^2+\rho^2).
\end{equation}
For $\rho>0$ the exact limiting density at zero is
\begin{equation}\label{eq:selection-coldens}
\hat p_{\rm col}(X,Y,Z;b,\kappa)=(2\pi)^{-3/2}\bigl[\rho^6(4X^2+\rho^2)\bigr]^{-1/2}\exp\Bigl[-\frac{b^2}4-\frac{\kappa^2(X^2-\frac14)^2}{\rho^2(4X^2+\rho^2)}\Bigr],
\end{equation}
and consequently
\begin{equation}\label{eq:selection-colscale}
p_{\nabla f(y)\mid\text{pair pins}}(0)=r^{-4}\bigl[\hat p_{\rm col}+o(1)\bigr]
\end{equation}
uniformly on compact subcharts away from the removable collision and axis strata. If $X$ stays away from the pair endpoints, the mean mismatch supplies exponential suppression at least $\exp(-c/\rho^2)$, and at $X=0$ it supplies $\exp(-c/\rho^4)$. Near $X=\pm\frac12$, lack of mean suppression occurs only inside a collision cone $|X\mp\frac12|=O(\rho)$.
\end{lemma}

\begin{proof}
The conditional moments \eqref{eq:selection-colmeans}--\eqref{eq:selection-Cperp} are exact Gaussian regressions computed from $\Cov(\partial^\alpha f,\partial^\beta f)=(-1)^{|\beta|}\partial^{\alpha+\beta}K_0(0)$, and \eqref{eq:selection-coldet} is the product of the axial variance \eqref{eq:selection-colvar} and $\det C_\perp$. The factor $r^{-4}$ in \eqref{eq:selection-colscale} is the Jacobian of the scaling \eqref{eq:selection-Gcol}. The suppression statements are read off from the exponent in \eqref{eq:selection-coldens}.
\end{proof}

\begin{lemma}[Hessian collision zero]\label{lem:selection-collision}
If $\nabla f(p)=\nabla f(q)=0$ and $f$ has bounded third derivatives between $p$ and $q$, then
\begin{equation}\label{eq:selection-collision}
\|H_p(q-p)\|\le C|q-p|^2,\qquad\|H_q(q-p)\|\le C|q-p|^2 .
\end{equation}
Hence each Hessian has a singular value $O(|q-p|)$, and its determinant has an additional factor $|q-p|$ times a polynomial bound in the remaining Hessian entries. The statement is dimension independent.
\end{lemma}

\begin{proof}
Subtract $\nabla f(q)-\nabla f(p)=0$ and use the integral formula $0=\int_0^1H_{p+s(q-p)}(q-p)\,ds$. Subtract $H_p(q-p)$ and use the Lipschitz continuity of $H$. The statement at $q$ is identical.
\end{proof}

\begin{lemma}[kernel-independent collar powers]\label{lem:selection-collarfactor}
Write a third point in the pair-scale collar as
\begin{equation}\label{eq:selection-collarpt2}
y=x+r(Xt+\rho v),
\end{equation}
where $v$ is a unit transverse direction, $|X|$ is bounded on the collar chart and $\rho\ge0$, and let $C_{\rm col}(\mathcal K;t,X,\rho,v)$ be the contact-limit conditional covariance of the scaled third gradient
\begin{equation}\label{eq:selection-colgrad}
\bigl(r^{-2}\partial_tf(y),\ r^{-1}\partial_vf(y),\ r^{-1}\partial_wf(y)\bigr),
\end{equation}
with $w$ the second transverse direction. For every $\mathcal K$ satisfying Hypothesis~\ref{hyp:selection-jet},
\begin{equation}\label{eq:selection-colfactor}
\det C_{\rm col}=\rho^6(X^2+\rho^2)\,A_{\rm col}(\mathcal K;t,X,\rho,v),
\end{equation}
where $A_{\rm col}$ extends continuously to $\rho=0$ and to $(X,\rho)=(0,0)$ after the ordinary radial blow-up, and the extension is analytic in $J_6(\mathcal K)$. At the planar jet,
\begin{equation}\label{eq:selection-Acol0}
A_{\rm col}(K_0)=\frac{4X^2+\rho^2}{X^2+\rho^2},
\end{equation}
whose radial-blow-up minimum is one. For every $\mathcal K$ satisfying Hypothesis~\ref{hyp:selection-jet} there is $a_{\rm col}(\mathcal K)>0$ such that
\begin{equation}\label{eq:selection-Acol}
A_{\rm col}(\mathcal K;t,X,\rho,v)\ge a_{\rm col}(\mathcal K)\qquad\text{for all }t,\ v\text{ and }(X,\rho)\ne(0,0).
\end{equation}
\end{lemma}

\begin{proof}
Take derivatives in the orthonormal frame $(t,v,w)$, so that $y=x+r(Xt+\rho v)$; in this frame the pins are the eight evaluations of \eqref{eq:setting-list}, up to an invertible linear map that does not depend on the kernel. Put
\[
U_1=f_{ttv},\qquad U_2=\tfrac12f_{tvv},\qquad W=f_{vvv},
\]
all at $x$. By \eqref{eq:selection-L1}--\eqref{eq:selection-Lperp}, read in this frame, the contact limit of the scaled gradient \eqref{eq:selection-colgrad} is
\[
\Bigl(\tfrac12\bigl(X^2-\tfrac14\bigr)f_{ttt}+\rho\,(XU_1+\rho U_2),\ \ \rho f_{vv},\ \ \rho f_{vw}\Bigr).
\]
This is an identity between Taylor coefficients, so it holds for every kernel. The term $f_{ttt}$ is a pin. Taking out the factor $\rho$ of each of the three entries gives
\[
\det C_{\rm col}=\rho^6D,\qquad D=\det\Cov\bigl(XU_1+\rho U_2,\ f_{vv},\ f_{vw}\bigm|P\bigr).
\]
Since $XU_1+\rho U_2$ is linear in $(X,\rho)$, $D$ is a quadratic form in $(X,\rho)$, and its coefficients are continuous in $(t,v)$ and analytic in $J_6(\mathcal K)$. Put $A_{\rm col}=D/(X^2+\rho^2)$. It is homogeneous of degree zero in $(X,\rho)$, so it is a continuous function of $(t,v)$ and of the direction of $(X,\rho)$; this is the radial blow-up, and \eqref{eq:selection-colfactor} holds.

For $(X,\rho)\ne(0,0)$ the polynomial symbols of the eleven variables $P$, $XU_1+\rho U_2$, $f_{vv}$, $f_{vw}$ are linearly independent, up to constant factors. The pins have the eight distinct monomials of \eqref{eq:setting-list}. The variable $XU_1+\rho U_2$ has the nonzero symbol $X\xi_t^2\xi_v+\frac\rho2\xi_t\xi_v^2$. The last two have the monomials $\xi_v^2$ and $\xi_v\xi_w$. None of these monomials occurs twice. By Hypothesis~\ref{hyp:selection-jet}(2) (for $K_{24}$, Theorem~\ref{thm:setting-gram}; for $K_0$, Remark~\ref{rem:selection-side24}), the joint covariance of these variables is positive definite. Hence so is its Schur complement given $P$, and $D>0$. Thus $A_{\rm col}$ is continuous and positive on a compact set: $t\in S^2$, $v$ a unit vector orthogonal to $t$, and the direction of $(X,\rho)$ with $\rho\ge0$. Its minimum $a_{\rm col}(\mathcal K)$ is therefore positive, which is \eqref{eq:selection-Acol}. At the planar jet, \eqref{eq:selection-coldet} gives \eqref{eq:selection-Acol0}: there $\Var(U_1\mid P)=2$, $\Var(U_2\mid P)=\frac12$, $\Cov(U_1,U_2\mid P)=0$, and $f_{vv}$, $f_{vw}$ are conditionally independent of $U_1,U_2$ with $\det\Cov(f_{vv},f_{vw}\mid P)=2$, so $D=4X^2+\rho^2$. The symbolic check \path{checks/s5_factor_check.py} verifies these identities.
\end{proof}

Equation \eqref{eq:selection-colfactor} is (3.14) of the main text.

\begin{lemma}[collar count]\label{lem:selection-collar}
For the side-$24$ field, uniformly for $t\in S^2$ and on compact positive mark sets,
\begin{equation}\label{eq:selection-colcount}
\E^{MS}_rN_{\rm col}\le C_{\rm col}\,r^3 .
\end{equation}
\end{lemma}

\begin{proof}
In a collar configuration with three critical points $M,S,y$, the points span a two-dimensional plane unless $\rho=0$. In the hard transverse direction orthogonal to that plane the Hessian block remains $O(1)$. The three station equations make the mixed hard/plane derivatives $O(r)$ after a fixed-chart divided-difference transformation. The Schur complement therefore reduces every full three-dimensional Hessian determinant to an $O(1)$ hard factor times the two-dimensional rank-two critical-contact determinant. Let $c$ be the cosine of the angle between $t$ and $y-x$, and fix $c_0\in(0,\frac14]$. Where $|c|\ge c_0$, the exact elimination of Lemma~\ref{lem:selection-elimination}, read in physical units, makes each two-dimensional determinant $r^2$ times the chart expression \eqref{eq:selection-B1}, \eqref{eq:selection-B1S} or \eqref{eq:selection-B1y}. Where $|c|<c_0$ those expressions are not used, since their denominators vanish at $c=0$. There $|X|\le\rho/2$, with $\rho$ the transverse coordinate of \eqref{eq:selection-collarpt}, and the sine of each angle of the triangle with vertices $M$, $S$ and $y$ is at least $c'\rho$, with $c'>0$ depending only on $C_1$. At each vertex $p$, Lemma~\ref{lem:selection-collision}, applied toward the other two vertices, gives unit vectors $u_1,u_2$ with $\sin\angle(u_1,u_2)\ge c'\rho$ and $\|H_pu_i\|\le Cr$. Since $|\det H|\le\|Hu_1\|\,\|Hu_2\|\,\|H\|/\sin\angle(u_1,u_2)$ for a $3\times3$ matrix $H$, this gives $|\det H_p|\le C\rho^{-1}r^2\Psi$ for $p\in\{M,S,y\}$, where $\Psi$ bounds the derivatives of orders two and three. The factor $\rho^{-3}$ is absorbed by the density: for small $\rho$ these points have $X$ near $0$, where Lemma~\ref{lem:selection-collarchart} gives the factor $\exp(-c/\rho^4)$. Either way, this gives a factor $r^2$ in each of $\det H_M$, $\det H_S$ and $\det H_y$, and therefore
\begin{equation}\label{eq:selection-colprod}
|\det H_M\det H_S\det H_y|\le Cr^6\,\hat D(X,Y,Z,U),
\end{equation}
where $U$ collects the random jet coefficients and $\hat D$ has uniformly bounded Gaussian polynomial moments. Near collision with $M$ or $S$, Lemma~\ref{lem:selection-collision} supplies an additional normalized squared-distance factor. This cancels the $\rho^{-3}$ density singularity in \eqref{eq:selection-coldens}: after integrating $X$ over the collision cone and using the transverse volume $\rho\,d\rho$, the local radial power is at least $\rho$. Away from collisions, the exponential axis suppression of Lemma~\ref{lem:selection-collarchart} is integrable. Collecting powers, $r^{-4}$ from the third-gradient density \eqref{eq:selection-colscale}, $r^6$ from \eqref{eq:selection-colprod}, $r^{-2}$ from the pair Palm normalizer \eqref{eq:selection-normalizer} and $r^3$ from the physical collar volume multiply to $r^3$. This proves \eqref{eq:selection-colcount} at the planar jet.

For the side-$24$ field, the Hessian collision zero, the two-dimensional plane elimination and the hard-complement Schur estimates are deterministic Taylor statements; they do not require the exact planar covariance. Under Hypothesis~\ref{hyp:selection-jet}, which $K_{24}$ satisfies by Remark~\ref{rem:selection-side24}, their random polynomial coefficients have uniformly bounded Gaussian moments, so the collar product bound $r^6$ times a Gaussian polynomial holds with changed constants only. By Lemmas~\ref{lem:selection-pinstab} and~\ref{lem:selection-collarfactor}, the collar zero-density estimate, the collision-cone integration and the $O(r^3)$ count hold for every kernel satisfying Hypothesis~\ref{hyp:selection-jet}, with constants depending on the kernel; in particular they hold for $K_{24}$. For such a kernel the axial and transverse blocks are again conditionally independent, because odd and even derivatives are uncorrelated for an even kernel. In the notation of the proof of Lemma~\ref{lem:selection-collarfactor}, the conditional axial mean differs from $\kappa(X^2-\frac14)$ by $\rho\,(X\,\E[U_1\mid P]+\rho\,\E[U_2\mid P])$. This is $O(\rho\kappa(|X|+\rho))$ on compact positive mark sets, since $U_1$ and $U_2$ are odd and the only odd pin with a nonzero target is $f_{ttt}$. Hence the suppression statements of Lemma~\ref{lem:selection-collarchart} persist for small $\rho$, with changed constants.
\end{proof}

\subsubsection{The singular-near count}

\begin{lemma}[singular-near corrected value--gradient frame at the planar jet]\label{lem:selection-singchart}
Let
\begin{equation}\label{eq:selection-singpt}
y=x+s\,\omega,\qquad\omega=(c,Y,Z),\qquad c^2+\rho^2=1,\qquad\eta=r/s,
\end{equation}
with $\rho^2=Y^2+Z^2$ and $\eta$ in a fixed compact interval $[0,\eta_0]$, $\eta_0<1$. Let the third value lie in the pair window,
\begin{equation}\label{eq:selection-window}
f(y)=b-\alpha\kappa r^3/6,\qquad\alpha\in[0,1],
\end{equation}
and use the corrected variables
\begin{gather}
V_0=s^{-3}\Bigl[f(y)-f(x)-\frac s2\bigl(Yf_2(y)+Zf_3(y)\bigr)-\frac{sc}3f_1(y)\Bigr],\label{eq:selection-V0}\\
V_1=s^{-2}f_1(y),\qquad V_2=s^{-1}f_2(y),\qquad V_3=s^{-1}f_3(y).\label{eq:selection-V123}
\end{gather}
The transformation from $(f(y),\nabla f(y))$ to $(V_0,V_1,V_2,V_3)$ has determinant of magnitude $s^{-7}$. At contact, with $\eta$ fixed, the random residual limits are
\begin{gather}
L_0=-\tfrac1{12}\bigl[Y\,f_{2ij}\omega_i\omega_j+Z\,f_{3ij}\omega_i\omega_j\bigr]-\tfrac{\eta^2}{16}\bigl[Yf_{112}+Zf_{113}\bigr],\label{eq:selection-L0}\\
L_1=\tfrac12f_{1ij}\omega_i\omega_j,\label{eq:selection-L1s}\\
(L_2,L_3)=Q(Y,Z)^{\top},\label{eq:selection-L23}
\end{gather}
with summation over repeated indices, and with the pinned term $\frac12c^2f_{111}$ removed from $L_1$. The other $\eta$-dependent terms are deterministic: $-\kappa\eta^2/4$ in $V_1$ and $-c\kappa\eta^2/6$ in $V_0$. The last term of $L_0$ is random. It arises because the pins force $f_k(x)=-\frac{r^2}8f_{k11}(x)+O(r^4)$ for $k=2,3$, and $V_0$ multiplies $f_2(x)$ and $f_3(x)$ by $Y/(2s^2)$ and $Z/(2s^2)$. The exact conditional covariance of $(L_0,L_1)$ is
\begin{equation}\label{eq:selection-C01}
C_{00}=\frac{\rho^2\bigl[(c^2+\frac34\eta^2)^2+4c^2\rho^2+3\rho^4\bigr]}{72},\qquad C_{01}=-\frac{c\rho^2(c^2+\rho^2+\frac34\eta^2)}6,\qquad C_{11}=\frac{\rho^2(4c^2+\rho^2)}2,
\end{equation}
so that
\begin{equation}\label{eq:selection-det01}
\det C_{(L_0,L_1)}=\frac{\rho^6Q_\eta}{144},\qquad Q_\eta=\bigl(3c^2-\tfrac34\eta^2\bigr)^2+12c^2\rho^2+3\rho^4 .
\end{equation}
Combining this with $\det C_\perp=2\rho^4$ gives
\begin{equation}\label{eq:selection-det4}
\det\Cov(L_0,L_1,L_2,L_3\mid P)=\frac{\rho^{10}Q_\eta}{72}.
\end{equation}
At $\eta=0$, $Q_0=3(c^2+\rho^2)(3c^2+\rho^2)$, so on the unit sphere
\begin{equation}\label{eq:selection-det4sphere}
\det\Cov(L_0,L_1,L_2,L_3\mid P)\big|_{\eta=0}=\frac{\rho^{10}(3c^2+\rho^2)}{24}.
\end{equation}
The conditional mean mismatch in $V_1$ is
\begin{equation}\label{eq:selection-mismatch}
\kappa\bigl(c^2-\eta^2/4\bigr).
\end{equation}
Consequently the angular density satisfies, near the axis,
\begin{equation}\label{eq:selection-pcorr}
p_{\rm corr}\le C\rho^{-5}\exp(-c_*/\rho^2)
\end{equation}
for a constant $c_*>0$, and, for every fixed $m\ge0$,
\begin{equation}\label{eq:selection-singdens}
\int_{S^2}\rho^{-m}\,p_{(f(y),\nabla f(y))\mid\text{pair pins}}(u,0)\,d\omega\le C_m\,s^{-7}.
\end{equation}
\end{lemma}

\begin{proof}
The transformation \eqref{eq:selection-V0}--\eqref{eq:selection-V123} is triangular with diagonal scalings $s^{-3},s^{-2},s^{-1},s^{-1}$, whence the determinant $s^{-7}$. The residual limits follow from Taylor's formula at $x$ and the pins, with $s=r/\eta$; \path{checks/s5_factor_check.py} verifies them, including the last term of $L_0$, for a polynomial field. The covariances \eqref{eq:selection-C01}--\eqref{eq:selection-det4} are exact Gaussian regressions computed from $\Cov(\partial^\alpha f,\partial^\beta f)=(-1)^{|\beta|}\partial^{\alpha+\beta}K_0(0)$, and \eqref{eq:selection-det4sphere} uses $c^2+\rho^2=1$. Near the pair axis, $|c|$ tends to one and $\eta\le\eta_0<1$, so \eqref{eq:selection-mismatch} is bounded away from zero. For any positive definite covariance, the Mahalanobis penalty is at least the squared mismatch of one coordinate divided by its variance. With the $V_1$ coordinate and $C_{11}\le2\rho^2$ on the unit sphere, the penalty is at least $c_0\kappa^2/\rho^2$. Also $Q_\eta\ge(3c^2-\frac34\eta^2)^2\ge(\frac94-3\rho^2)^2$ for $\rho^2<\frac34$, so the angular density has the uniform form \eqref{eq:selection-pcorr} near the axis, which is integrable against the $S^2$ area element $\rho\,d\rho\,d\phi$. Away from the axis, $Q_\eta\ge3\rho^4$ gives a positive compact eigenfloor, uniformly in $\eta\in[0,\eta_0]$. Thus the exact finite-$r$ corrected covariance extends continuously and, with the Jacobian $s^{-7}$, supplies \eqref{eq:selection-singdens}; the weight $\rho^{-m}$ is integrable against \eqref{eq:selection-pcorr} near the axis, because $\exp(-c_*/\rho^2)$ decays faster than any power of $\rho$.
\end{proof}

The bound \eqref{eq:selection-singdens} is the density half of the singular-near witness bound; the determinant half is Lemma~\ref{lem:selection-envelope}.

\begin{lemma}[kernel-independent extra $\rho^2$ cancellation]\label{lem:selection-singfactor}
Write $y=x+s\omega$ with $\omega=(c,\rho v)$, $c^2+\rho^2=1$ and $\eta=r/s\in[0,\eta_0]$, $\eta_0<1$. Use the corrected variables $V_0,\dots,V_3$ of Lemma~\ref{lem:selection-singchart}, and let $C_{\rm sing}$ denote their contact-limit conditional covariance. For every $\mathcal K$ satisfying Hypothesis~\ref{hyp:selection-jet},
\begin{equation}\label{eq:selection-singfactor}
\det C_{\rm sing}=\rho^{10}(3c^2+\rho^2)\,A_{\rm sing}(\mathcal K;t,c,\rho,v,\eta),
\end{equation}
where $A_{\rm sing}$ extends continuously to $\rho=0$ and is analytic in $J_6(\mathcal K)$. At the planar jet,
\begin{equation}\label{eq:selection-Asing0}
A_{\rm sing}(K_0)=\frac{Q_\eta}{72(3c^2+\rho^2)},\qquad\text{which is }\frac1{24}\text{ at }\eta=0,
\end{equation}
with $Q_\eta$ as in \eqref{eq:selection-det01}. For every $\mathcal K$ satisfying Hypothesis~\ref{hyp:selection-jet} there is $a_{\rm sing}(\mathcal K)>0$ such that
\begin{equation}\label{eq:selection-Asing}
A_{\rm sing}(\mathcal K;t,c,\rho,v,\eta)\ge a_{\rm sing}(\mathcal K)\qquad\text{for all }t,\ v,\ c^2+\rho^2=1\text{ and }\eta\in[0,\eta_0].
\end{equation}
Thus the $\rho^{-5}$ density factor and the $\exp(-c_*/\rho^2)$ axis suppression persist for small $\rho$, with changed constants. The conditional mean mismatch in $V_1$ is $\kappa(c^2-\eta^2/4)+O(\rho\kappa)$, which is separated from zero in the axis chart because $\eta_0<1$.
\end{lemma}

\begin{proof}
Take derivatives in the frame $(t,v,w)$, so that $\omega=ct+\rho v$, and let $U_1$, $U_2$ and $W$ be as in the proof of Lemma~\ref{lem:selection-collarfactor}. By \eqref{eq:selection-L0}--\eqref{eq:selection-L23}, read in this frame, the random residuals are, for every kernel,
\begin{equation}\label{eq:selection-singlead}
L_0=-\tfrac{\rho}{12}\bigl((c^2+\tfrac34\eta^2)U_1+4c\rho\,U_2+\rho^2W\bigr),\qquad L_1=\rho\,(cU_1+\rho U_2),\qquad (L_2,L_3)=\rho\,(f_{vv},f_{vw}).
\end{equation}
Thus $(L_0,L_1)=(-\frac\rho{12}\tilde L_0,\ \rho\tilde L_1)$ with $(\tilde L_0,\tilde L_1)^{\top}=A\,(U_1,U_2,W)^{\top}$ and
\[
A=\begin{pmatrix}c^2+\frac34\eta^2&4c\rho&\rho^2\\ c&\rho&0\end{pmatrix}.
\]
Every $2\times2$ minor of $A$ is divisible by $\rho$: the minors are $\rho\,\tilde m$ with $\tilde m=(\frac34\eta^2-3c^2,\,-c\rho,\,-\rho^2)$. Let $\Sigma=\Cov(U_1,U_2,W\mid P)$. By the Cauchy--Binet formula,
\[
\det\Cov(\tilde L_0,\tilde L_1\mid P)=\det(A\Sigma A^{\top})=\rho^2Q_\eta(\mathcal K),\qquad Q_\eta(\mathcal K)=\sum_{S,T}\tilde m_S\,\tilde m_T\det\Sigma_{S,T},
\]
the sum running over pairs of two-element column sets. For an even kernel, odd and even derivatives are uncorrelated, so $(L_0,L_1)$ and $(L_2,L_3)$ are conditionally independent given $P$. Hence
\[
\det C_{\rm sing}=\frac{\rho^{10}}{144}\,Q_\eta(\mathcal K)\,\det\Cov(f_{vv},f_{vw}\mid P).
\]
The nominal power $\rho^8$ is thus $\rho^{10}$; this is the cancellation, and it holds for every even kernel. Put $A_{\rm sing}=Q_\eta(\mathcal K)\det\Cov(f_{vv},f_{vw}\mid P)/\bigl(144(3c^2+\rho^2)\bigr)$. It is continuous in $(t,v,c,\rho,\eta)$ and analytic in $J_6(\mathcal K)$, and \eqref{eq:selection-singfactor} holds.

The symbols of $P$, $U_1$, $U_2$, $W$, $f_{vv}$ and $f_{vw}$ are distinct monomials. As in the proof of Lemma~\ref{lem:selection-collarfactor}, $\Sigma$ and $\Cov(f_{vv},f_{vw}\mid P)$ are therefore positive definite. Since $A\Sigma A^{\top}\succcurlyeq\lambda_{\min}(\Sigma)AA^{\top}$ and $\det(AA^{\top})=\rho^2|\tilde m|^2$, we get $Q_\eta(\mathcal K)\ge\lambda_{\min}(\Sigma)^2|\tilde m|^2$. Moreover $\tilde m\ne0$: if $\rho=0$ then $c^2=1$, and $\frac34\eta^2-3\ne0$ because $\eta<1$. So $A_{\rm sing}$ is continuous and positive on the compact set given by $t\in S^2$, $v$ a unit vector orthogonal to $t$, $c^2+\rho^2=1$, $\rho\ge0$ and $\eta\in[0,\eta_0]$. Its minimum $a_{\rm sing}(\mathcal K)$ is positive, which is \eqref{eq:selection-Asing}. At the planar jet $\Sigma=\operatorname{diag}(2,\frac12,6)$ and $\det\Cov(f_{vv},f_{vw}\mid P)=2$, so $Q_\eta(K_0)$ is the $Q_\eta$ of \eqref{eq:selection-det01}, and \eqref{eq:selection-det4} gives \eqref{eq:selection-Asing0}. The check \path{checks/s5_factor_check.py} verifies these identities.

The means. The random parts in \eqref{eq:selection-singlead} are odd, and the only odd pin with a nonzero target is $f_{ttt}$. So the conditional mean of $V_1$ differs from its planar value $\kappa(c^2-\eta^2/4)$ by the difference of two regressions of $\rho(cU_1+\rho U_2)$ on the pins, which is $O(\rho\kappa)$ on compact positive mark sets. For $\rho\le\rho_0(\mathcal K)$ the mismatch is therefore separated from zero. As in the proof of Lemma~\ref{lem:selection-singchart}, the Mahalanobis penalty is at least the squared $V_1$ mismatch over $C_{11}=\Var(L_1\mid P)\le C\rho^2$, hence at least $c_*\kappa^2/\rho^2$.
\end{proof}

Equation \eqref{eq:selection-singfactor} is (3.15) of the main text.

\begin{lemma}[anisotropic triple-determinant envelope]\label{lem:selection-envelope}
Let $y=x+s\omega$ be a critical point of $f$ in the singular-near region, $C_1r<s\le\delta$, with $\omega=ct+\rho n$, where $n\perp t$ is a unit vector, $c^2+\rho^2=1$ and $\rho>0$, and let $f(y)$ lie in the pair window \eqref{eq:selection-window}. Let $\Psi\ge1$ be a fixed polynomial in $\kappa$ and in the suprema over the ball $B_{3\delta}(x)$ of the norms of the derivatives of $f$ of orders two to five. Under the pair pins,
\begin{gather}
|\det H_M|+|\det H_S|\le C\rho^{-2}\,r(r+s^2)\,\Psi,\label{eq:selection-endenv}\\
|\det H_y|\le C\rho^{-2}\,(r^2+s^3)\,\Psi,\label{eq:selection-thirdenv}
\end{gather}
where $C$ depends only on $C_1$. Consequently
\begin{equation}\label{eq:selection-envelope}
\E\bigl[W_r|\det H_y|\bigm|\text{pins},\ f(y)=u,\ \nabla f(y)=0\bigr]\le C\rho^{-6}\,\bar\Psi(\rho)\,[r(r+s^2)]^2(r^2+s^3),
\end{equation}
where $\bar\Psi(\rho)$, the conditional moment of $\Psi^3$, is at most $C(1+\rho^{-m_0})$ for a fixed $m_0$, uniformly on compact marks and over the window. The same bounds hold for every kernel satisfying Hypothesis~\ref{hyp:selection-jet}, in particular for $K_{24}$, with constants depending on the kernel.
\end{lemma}

\begin{proof}
All estimates are deterministic consequences of Taylor's formula on $B_{3\delta}(x)$, which contains every segment used below. We write $O(X)$ for a quantity bounded by $C\Psi X$. Since $C_1>1$ and $r\le\delta$, the distances from $y$ to $M$ and to $S$ lie in $[s/2,2s]$, and $r<s$; these inequalities absorb the lower-order terms below. The elimination of Lemma~\ref{lem:selection-elimination} is not used.

\emph{Expansions.} For a critical point $p$ of $f$ put $A=H_p$, $T=D^3f(p)$, $\Theta=D^4f(p)$ and $F_p(v)=f(p+v)-f(p)-\frac12\nabla f(p+v)\cdot v$. For $|v|\le2s$,
\begin{gather}
\nabla f(p+v)=Av+\tfrac12T[v,v]+\tfrac16\Theta[v,v,v]+R_p(v),\label{eq:selection-taylorgrad}\\
F_p(v)=-\tfrac1{12}T[v,v,v]-\tfrac1{24}\Theta[v,v,v,v]+R'_p(v),\label{eq:selection-taylorF}
\end{gather}
with $|R_p(v)|+|v|\,|DR_p(v)|\le C\Psi|v|^4$ and $|R'_p(v)|+|v|\,|DR'_p(v)|\le C\Psi|v|^5$; the coefficients of $F_p$ are $\frac1{k!}(1-\frac k2)$, which vanish for $k=2$. If $p+v$ is also a critical point, the left side of \eqref{eq:selection-taylorgrad} vanishes and $F_p(v)=f(p+v)-f(p)$.

\emph{The endpoint $M$.} Take $p=M$. Then $S=M+rt$ and $y=M+w$ with $w=\frac r2t+s\omega$ and $|w|\in[s/2,2s]$; put $\hat w=w/|w|$. The gradient equation at $S$ gives $At=-\frac r2T[t,t]+O(r^2)$, so $A[t,e]=O(r)$ for every unit vector $e$. The gradient equation at $y$ gives $Aw=-\frac12T[w,w]+O(s^3)$, so $A[\hat w,e]=O(s)$. Dotting it with $w$ and eliminating $T[w,w,w]$ with \eqref{eq:selection-taylorF} gives
\[
A[w,w]=6F_M(w)+\tfrac1{12}\Theta[w,w,w,w]+O(s^5),
\]
and the window gives $F_M(w)=f(y)-f(M)\in[-\kappa r^3/6,0]$; hence $A[\hat w,\hat w]=O(r^3s^{-2}+s^2)$. Let $e_h$ be a unit vector orthogonal to $t$ and $\omega$. In the basis $(t,\hat w,e_h)$ the entries of $A$ are $O(r)$ in the first row and column, $A[\hat w,\hat w]=O(r^3s^{-2}+s^2)$, $A[\hat w,e_h]=O(s)$ and $A[e_h,e_h]=O(1)$. Expanding along the first row, every term of $\det\bigl(A[u_i,u_j]\bigr)$ is $O(r(r+s^2))$. That determinant equals $\det H_M$ times the Gram determinant $\sin^2\angle(t,\hat w)$, and $\sin\angle(t,\hat w)=s\rho/|w|\ge\rho/2$. This gives \eqref{eq:selection-endenv} at $M$. The endpoint $S$ is the same with $t$ replaced by $-t$, using the window in the form $f(y)-f(S)\in[0,\kappa r^3/6]$.

\emph{The third point.} Take $p=y$, $a=M-y$ and $b=S-y$, so that $b-a=rt$ and $\frac12(a+b)=x-y=-s\omega$. Subtracting and adding the gradient equations at $a$ and $b$, with $|R_y(b)-R_y(a)|\le C\Psi rs^3$, gives
\begin{gather}
At=s\,T[\omega,t]-\tfrac{s^2}2\Theta[\omega,\omega,t]+O(r^2+s^3),\label{eq:selection-At}\\
A\omega=\tfrac s2T[\omega,\omega]+\tfrac{r^2}{8s}T[t,t]-\tfrac{s^2}6\Theta[\omega,\omega,\omega]+O(r^2+s^3).\label{eq:selection-Aw}
\end{gather}
The pins give $F_y(a)-F_y(b)=f(M)-f(S)=\kappa r^3/6$, and the window gives $0\le F_y(a)\le\kappa r^3/6$. Expanding \eqref{eq:selection-taylorF} at $a$ and $b$, with $|R'_y(a)-R'_y(b)|\le C\Psi rs^4$, gives
\begin{gather}
T[\omega,\omega,t]=\tfrac23s\,\Theta[\omega,\omega,\omega,t]+O(r^2s^{-2}+s^2),\label{eq:selection-Twwt}\\
T[\omega,\omega,\omega]=\tfrac12s\,\Theta[\omega,\omega,\omega,\omega]+O(r^2s^{-2}+s^2);\label{eq:selection-Twww}
\end{gather}
in \eqref{eq:selection-Twww}, \eqref{eq:selection-Twwt} cancels the term $\frac1{12}s^3r\,\Theta[\omega,\omega,\omega,t]$. Inserting these into \eqref{eq:selection-At}--\eqref{eq:selection-Aw},
\begin{gather*}
A[t,t]=O(s),\qquad A[\omega,t]=\tfrac16s^2\Theta[\omega,\omega,\omega,t]+O(r^2s^{-1}+s^3),\\
A[\omega,\omega]=\tfrac1{12}s^2\Theta[\omega,\omega,\omega,\omega]+O(r^2s^{-1}+s^3),
\end{gather*}
so $A[t,t]A[\omega,\omega]-A[\omega,t]^2=O(r^2+s^3)$. The entries $A[t,e_h]$ and $A[\omega,e_h]$ are $O(s)$ by \eqref{eq:selection-At}--\eqref{eq:selection-Aw}, and $A[e_h,e_h]=O(1)$. Expanding the $3\times3$ determinant in the basis $(t,\omega,e_h)$ along the last row, every term is $O(r^2+s^3)$. The Gram determinant of this basis is $\rho^2$, which gives \eqref{eq:selection-thirdenv}.

\emph{Conditional moments.} Under the conditioning on the pins and on $(f(y),\nabla f(y))=(u,0)$, every derivative field of $f$ is Gaussian, with mean given by regression on the corrected variables of Lemmas~\ref{lem:selection-pairlaw}--\ref{lem:selection-singchart} and covariance dominated by the unconditional one. The corrected pair frame has a uniform eigenfloor, and by \eqref{eq:selection-det4} and Lemma~\ref{lem:selection-singfactor} the corrected third-point block has determinant at least a constant times $\rho^{10}$. So the regression coefficients, and with them the conditional means, are at most $C(1+\rho^{-m})$ on compact marks and over the window, for a fixed $m$. By the Borell--TIS inequality \cite{AdlerTaylor}, every conditional moment of $\Psi$ is at most a constant times a power of $1+\rho^{-1}$. Multiplying the three envelopes, with $W_r\le|\det H_M|\,|\det H_S|$, gives \eqref{eq:selection-envelope}.

The argument uses only Taylor's formula, Gaussian moments and the lower bound $\rho^{10}$ for the third-point block, which Lemma~\ref{lem:selection-singfactor} gives for every kernel satisfying Hypothesis~\ref{hyp:selection-jet}. So it holds for every such kernel, with constants depending on the kernel.
\end{proof}

The envelope \eqref{eq:selection-envelope} is (3.16) of the main text; the factor $\rho^{-6}\bar\Psi(\rho)$ is absorbed by the angular integral \eqref{eq:selection-singdens}. The symbolic check \path{checks/s5_envelope_check.py} confirms the cancellations behind \eqref{eq:selection-endenv} and \eqref{eq:selection-thirdenv} for polynomial fields of degree seven: in each of the four determinants (the in-plane part and the hard-direction part, at an endpoint and at the third point), every monomial in $r$ and $s$ is within the stated envelope.

\begin{lemma}[singular-near count]\label{lem:selection-sing}
For the side-$24$ field, uniformly for $t\in S^2$ and on compact positive mark sets,
\begin{equation}\label{eq:selection-singcount}
\E^{MS}_rN_{\rm sing}\le C_{\rm sing}\,r^3 .
\end{equation}
\end{lemma}

\begin{proof}
By the Kac--Rice formula, $\E^0[W_rN_{\rm sing}]$ is the integral over $y$ in region~(ii) and over $u$ in the window of the conditional expectation in \eqref{eq:selection-envelope} times the conditional density of $(f(y),\nabla f(y))$ at $(u,0)$. Combine the envelope \eqref{eq:selection-envelope} with the density bound \eqref{eq:selection-singdens} for $m=6+m_0$, the pair normalizer $Z_r\ge c_Zr^2$ of \eqref{eq:selection-normalizer}, the value-window width $O(r^3)$ of \eqref{eq:selection-window} and the polar volume $s^2\,ds$. This gives
\begin{equation}\label{eq:selection-singint}
\E^{MS}_rN_{\rm sing}\le C\,r\int_{C_1r}^{\delta}s^{-5}\,[r(r+s^2)]^2(r^2+s^3)\,ds .
\end{equation}
Expanding the polynomial produces the six terms
\begin{equation}\label{eq:selection-sixterms}
r^7s^{-5},\qquad r^6s^{-3},\qquad r^5s^{-2},\qquad r^5s^{-1},\qquad r^4,\qquad r^3s^2,
\end{equation}
up to fixed numerical factors. Integration from $C_1r$ to $\delta$ bounds them respectively by $O(r^3)$, $O(r^4)$, $O(r^4)$, $O(r^5\log(1/r))$, $O(r^4)$ and $O(r^3)$. This proves \eqref{eq:selection-singcount} at the planar jet. For the side-$24$ field, Lemma~\ref{lem:selection-singfactor} shows that the $\rho^{-5}$ density factor, the $\exp(-c_*/\rho^2)$ axis suppression and the separated mean mismatch persist, so \eqref{eq:selection-singdens} holds with changed constants; with Lemma~\ref{lem:selection-envelope} the same power count applies.
\end{proof}

\subsubsection{The fixed-distance count}

\begin{lemma}[conditional density block]\label{lem:selection-farblock}
Let
\begin{equation}\label{eq:selection-farblock}
\mathcal A_r=\bigl(P^{\rm corr}_r,f(y),\nabla f(y)\bigr),\qquad\dist(y,x)>\delta,
\end{equation}
where $P^{\rm corr}_r$ is the corrected pair pin vector; then $\dist(y,\{M,S\})\ge\delta/2$. For the side-$24$ field, complete positive Fourier support and derivative-distribution independence give, uniformly on the compact separated configuration set,
\begin{equation}\label{eq:selection-farfloor}
c_\delta I\le\Sigma_{\mathcal A\mathcal A}\le C_\delta I .
\end{equation}
Let $L_r(b,\kappa)$ be the corrected pair-pin target and let the remote critical value be
\begin{equation}\label{eq:selection-farwindow}
u=b-\alpha\kappa r^3/6,\qquad\alpha\in[0,1].
\end{equation}
The full density target is
\begin{equation}\label{eq:selection-fartarget}
a_r(b,\kappa,\alpha)=\bigl(L_r(b,\kappa),u,0\bigr).
\end{equation}
On compact positive mark sets, the conditional density of the remote value and gradient given the pins satisfies
\begin{equation}\label{eq:selection-fardens}
p_{(f(y),\nabla f(y))\mid\text{pair pins}}(u,0)=\frac{p_{\mathcal A}(a_r)}{p_{P^{\rm corr}_r}(L_r(b,\kappa))}\le C .
\end{equation}
\end{lemma}

\begin{proof}
The contact-limit corrected pair frame consists of distinct derivative coordinates at $x$, and $(f(y),\nabla f(y))$ of derivative coordinates at $y$, so the associated distributions are linearly independent; full-spectrum Gram positivity (Section~\ref{sec:setting}) and compactness of the separated configuration set give \eqref{eq:selection-farfloor}. The pins $\Pi_{r,x,t}$ and $P^{\rm corr}_r$ are related by an invertible linear map, whose Jacobian cancels between numerator and denominator, so the conditional density is the ratio in \eqref{eq:selection-fardens}. By \eqref{eq:selection-farfloor} the numerator is at most $(2\pi)^{-6}c_\delta^{-6}$. On compact marks the target $L_r(b,\kappa)$ is bounded, and the covariance of $P^{\rm corr}_r$ has uniform eigenvalue bounds (Lemma~\ref{lem:selection-pinstab}), so the denominator is bounded below.
\end{proof}

\begin{lemma}[fixed-distance count]\label{lem:selection-far}
Let $\mathcal Y_\delta$ be the compact side-$24$ configuration region with $\dist(y,x)>\delta$, and let $H_y$ be the Hessian at $y$. On compact positive mark sets,
\begin{equation}\label{eq:selection-threedet}
\E\bigl[|\det H_M\det H_S\det H_y|\bigm|\mathcal A_r=a_r\bigr]\le Cr^2,
\end{equation}
and, uniformly for $t\in S^2$,
\begin{equation}\label{eq:selection-farcount}
\E^{MS}_rN_{\rm far}\le C_{\rm far}\,r^3 .
\end{equation}
\end{lemma}

\begin{proof}
Condition on $\mathcal A_r=a_r(b,\kappa,\alpha)$. All raw cross-covariances with the scaled endpoint variables of Lemma~\ref{lem:selection-endpoint} and with $H_y$ are bounded, and only $\Sigma_{\mathcal A\mathcal A}$ is inverted. Gaussian regression therefore gives a conditional mean bounded by $C(1+|b|+\kappa)$ and a bounded positive-semidefinite conditional covariance, even if the enlarged Hessian block is degenerate. Since the product of three $3\times3$ determinants has total degree nine, \eqref{eq:selection-softdet} and Lemma~\ref{lem:selection-moments} give \eqref{eq:selection-threedet} on compact marks.

By the Kac--Rice formula, $\E^0[W_rN_{\rm far}]$ is the integral over $y\in\mathcal Y_\delta$ and over $u\in[b-\kappa r^3/6,b]$ of the conditional expectation in \eqref{eq:selection-threedet} times the conditional density \eqref{eq:selection-fardens}. Hence
\begin{equation}\label{eq:selection-farnum}
\E^0[W_rN_{\rm far}]\le C\,r^2\cdot\frac{\kappa r^3}6\cdot|\mathcal Y_\delta|\le Cr^5 .
\end{equation}
Dividing explicitly by $Z_r=\E^0[W_r]\ge c_Zr^2$ from \eqref{eq:selection-normalizer} gives \eqref{eq:selection-farcount}.
\end{proof}

The bound \eqref{eq:selection-farcount} is (3.17) of the main text. The count is conditional on the pins, so the argument uses the lower bound on the pin density in \eqref{eq:selection-fardens}; it gives no Gaussian factor in $(b,\kappa)$, and it is stated only on compact marks, which is all that Theorem~\ref{thm:selection-main} needs.

\subsection{Capture and escape}\label{ssec:selection-escape}

\paragraph{Normalized potential.} On a fixed compact convex chart $D^+$ consider, for $X\in\R$ and $Y\in\R^2$,
\begin{equation}\label{eq:selection-potential}
P(X,Y)=p_0(X)+\tfrac12\bigl(X^2-\tfrac14\bigr)a\cdot Y-\tfrac12Y^{\top}TY+\tfrac12X\,Y^{\top}WY+\tfrac16C[Y,Y,Y]+r\,Q(X,Y),
\end{equation}
where $p_0(X)=X^3/3-X/4-1/12$, $T\succ0$ is symmetric, $W$ is symmetric, $C$ is a symmetric cubic tensor, and $Q(M)=Q(S)=0$, $\nabla Q(M)=\nabla Q(S)=0$ at $M=(-\frac12,0)$ and $S=(\frac12,0)$; note $P(M)=p_0(-\frac12)=0$. Put
\begin{equation}\label{eq:selection-R}
R=1+|a|+\|W\|+\|C\|+\|Q\|_{C^2(D^+)}\ge1,\qquad\tau=\lambda_{\min}(T).
\end{equation}
For the conditioned side-$24$ field, $T=-Q_\perp/(\kappa r)$ (see \eqref{eq:selection-Tr} below; the factor $1/r$ is essential). Fix
\begin{equation}\label{eq:selection-assume}
K=4096,\qquad\varepsilon_0=\frac1{1024K},\qquad\text{and assume}\qquad rR^5\le\varepsilon_0,\qquad\tau\ge2KR^5+2R .
\end{equation}
Write $E(Y)=Y^{\top}TY$ and $q(Y)=|TY|$; then $q^2\ge\tau E$ and $|Y|^2\le E/\tau$. The gradient of $P$ is
\begin{equation}\label{eq:selection-gradient}
\begin{aligned}
F_X&=X^2-\tfrac14+Xa\cdot Y+\tfrac12Y^{\top}WY+rQ_X,\\
F_Y&=\tfrac12\bigl(X^2-\tfrac14\bigr)a-TY+XWY+\tfrac12C[Y,Y,\cdot\,]+r\nabla_YQ,
\end{aligned}
\end{equation}
and the half-branches below are trajectories of the gradient-ascent flow $(\dot X,\dot Y)=(F_X,F_Y)$.

\begin{lemma}[endpoint energy slope]\label{lem:selection-slope}
The Jacobian at $S$ has block form $J_S=\begin{pmatrix}A&B^{\top}\\B&-T+E_D\end{pmatrix}$ with $A\ge3/4$, $|B|\le R$ and $\|E_D\|\le3R/4$ (after one fixed small-$r$ reduction), and exactly one positive eigenvalue. Its positive eigenvector $v=(v_X,v_Y)$ satisfies
\begin{equation}\label{eq:selection-slope}
E(v_Y)\le4\,(R^2/\tau)\,|v_X|^2 .
\end{equation}
\end{lemma}

\begin{proof}
From $J_Sv=\lambda_+v$, the transverse block gives $(\lambda_+I+T-E_D)v_Y=Bv_X$. Taking the inner product with $v_Y$,
\[
E(v_Y)=v_Y^{\top}Tv_Y\le|B|\,|v_X|\,|v_Y|+\|E_D\|\,|v_Y|^2\le R\,|v_X|\,|v_Y|+\tfrac{3R}4|v_Y|^2,
\]
using $\lambda_+>0$ and $v_Y^{\top}E_Dv_Y\le\|E_D\|\,|v_Y|^2$. Young's inequality $R|v_X||v_Y|\le\frac{R^2}\tau|v_X|^2+\frac\tau4|v_Y|^2$ and the energy bound $|v_Y|^2\le E(v_Y)/\tau$ give
\[
E(v_Y)\le\frac{R^2}\tau|v_X|^2+\Bigl(\frac\tau4+\frac{3R}4\Bigr)|v_Y|^2\le\frac{R^2}\tau|v_X|^2+\Bigl(\frac14+\frac38\Bigr)E(v_Y),
\]
where $\tau\ge2R$, from \eqref{eq:selection-assume}, gives $\frac{3R}4\cdot\frac1\tau\le\frac38$. Hence $\frac38E(v_Y)\le\frac{R^2}\tau|v_X|^2$, that is, $E(v_Y)\le\frac83(R^2/\tau)|v_X|^2\le4(R^2/\tau)|v_X|^2$. The strict final inequality leaves a fixed normalized margin, which persists under the exact $C^1$ perturbation of Lemma~\ref{lem:selection-robust}.
\end{proof}

\begin{lemma}[outward energy cone]\label{lem:selection-cone}
Let $\xi=X-\frac12$, $\delta_R=1/(512R^3)$ and
\[
\mathcal C_E=\bigl\{0\le\xi\le\delta_R,\ E(Y)\le4(R^2/\tau)\xi^2\bigr\}.
\]
On $\mathcal C_E$: $q\ge2R\xi$ on the nonzero boundary, $|Y|\le\xi/K$, $F_X\ge3\xi/4$, and $G_E=E(Y)-4(R^2/\tau)\xi^2$ satisfies
\begin{equation}\label{eq:selection-coneflux}
\tfrac12\dot G_E\le-\tfrac23q^2<0\qquad\text{on }\{G_E=0\}.
\end{equation}
Hence the outward half-branch stays in $\mathcal C_E$ until $\xi=\delta_R$.
\end{lemma}

\begin{proof}
On $G_E=0$ we have $E=4(R^2/\tau)\xi^2$, so $q\ge\sqrt{\tau E}=2R\xi$ and $|Y|\le\sqrt{E/\tau}=2R\xi/\tau\le\xi/(KR^4)$. (The full strength $\tau\ge2KR^5$ is used here, not merely $\tau\ge2KR$; this matters for the next bound when $R$ is large.) Then $|Xa\cdot Y|\le(\frac12+\delta_R)R|Y|\le(\frac12+\delta_R)\xi/(KR^3)\le\xi/16$, as $(\frac12+\delta_R)/(KR^3)\le1/K\le1/16$; next $\frac12|Y^{\top}WY|\le\frac R2|Y|^2\le\frac R2\cdot\frac{\xi^2}{K^2R^8}\le\frac\xi{16}\cdot\frac{8\xi}{K^2R^7}\le\xi/16$, as $8\xi\le8\delta_R\le1\le K^2R^7$; and, since $\nabla Q(S)=0$ and $\|Q\|_{C^2}\le R$, $|Q_X|\le R(\xi+|Y|)\le R(1+1/K)\xi$ on $\mathcal C_E$, so $|rQ_X|\le2rR\xi\le\xi/16$, as $2rR\le2\varepsilon_0R^{-4}\le2\varepsilon_0<1/16$. Since $X^2-\frac14=\xi(1+\xi)\ge\xi$, we get $F_X\ge\xi-3\xi/16\ge3\xi/4$.

For the flux, $\frac12\dot G_E=(TY)\cdot F_Y-4(R^2/\tau)\xi F_X$, and the moving-boundary term is nonpositive because $F_X\ge0$. In $(TY)\cdot F_Y$ the leading term is $-q^2$. We divide the four adverse terms by $q^2$, using $q\ge2R\xi$, $|Y|\le2R\xi/\tau$, $|X|\le\frac12+\delta_R$, $X^2-\frac14=\xi(1+\xi)$, and, for the last term, that $\nabla Q$ vanishes at $S$ by \eqref{eq:selection-potential}, so that $|\nabla_YQ|\le R(\xi+|Y|)$ on $\mathcal C_E$ by the $C^2$ bound:
\begin{equation}\label{eq:selection-coneledger}
\begin{aligned}
\frac{|(TY)\cdot\frac12(X^2-\frac14)a|}{q^2}&\le\frac{\frac12\xi(1+\xi)Rq}{q^2}=\frac{\xi(1+\xi)R}{2q}\le\frac{1+\xi}4\le\frac{1+\delta_R}4\le\frac{257}{1024},\\
\frac{|(TY)\cdot XWY|}{q^2}&\le\frac{|X|\,R\,|Y|}q\le\frac{(\frac12+\delta_R)R\,(2R\xi/\tau)}{2R\xi}=\frac{(\frac12+\delta_R)R}\tau\le\frac{257}{1024K},\\
\frac{|(TY)\cdot\frac12C[Y,Y,\cdot\,]|}{q^2}&\le\frac{\frac12R|Y|^2q}{q^2}\le\frac{R(2R\xi/\tau)^2}{2\cdot2R\xi}=\frac{R^2\xi}{\tau^2}\le\frac{\delta_R}{4K^2R^8}\le\frac1{2048K^2},\\
\frac{|(TY)\cdot r\nabla_YQ|}{q^2}&\le\frac{rR(\xi+|Y|)}q\le\frac{rR\xi(1+2R/\tau)}{2R\xi}=\frac{r(1+2R/\tau)}2\le2\varepsilon_0 .
\end{aligned}
\end{equation}
Here the second line uses $q\ge2R\xi$ and $|Y|\le2R\xi/\tau$, the third uses $\tau\ge2KR^5$ and $\xi\le\delta_R=1/(512R^3)$, and the fourth uses $r\le\varepsilon_0/R^5\le\varepsilon_0$ and $2R/\tau\le1/(KR^4)\le1$. The four brackets $\frac{257}{1024}$, $\frac{257}{1024K}$, $\frac1{2048K^2}$ and $2\varepsilon_0$ sum to $\frac{8625610753}{34359738368}<\frac13$ (exact rational arithmetic). Hence $\frac12\dot G_E\le-(1-\frac13)q^2$.
\end{proof}

\begin{lemma}[outward energy strip and level bound]\label{lem:selection-strip}
Let
\[
\mathcal S_E=\bigl\{\tfrac12+\delta_R\le X\le\tfrac54,\ E(Y)\le E_*\bigr\},\qquad E_*=\tfrac9{16}(R^2/\tau).
\]
On its lateral boundary, $q\ge3R/4$, $|Y|\le3R/(4\tau)$, $F_X\ge3\delta_R/4>0$, and
\begin{equation}\label{eq:selection-stripflux}
\tfrac12\dot E\le-\tfrac1{10}q^2<0\qquad\text{on }\{E=E_*\}.
\end{equation}
The branch therefore reaches $X=\frac54$ with
\begin{equation}\label{eq:selection-level}
P\bigl(\tfrac54,Y\bigr)>\frac{49}{192}-\frac1{768}=\frac{65}{256}>\frac14>0=P(M):
\end{equation}
it reaches a point strictly above the maximum level and cannot terminate at $M$.
\end{lemma}

\begin{proof}
At $\xi=\delta_R$, $E\le4(R^2/\tau)\delta_R^2=\frac{4\delta_R^2}{9/16}E_*<E_*$, since $\delta_R=1/(512R^3)<3/8$. On $\{E=E_*\}$, $q\ge\sqrt{\tau E_*}=3R/4$, and the adverse-to-$q^2$ ratios are $\frac78$, $\frac5{8K}$, $\frac3{32K^2}$ and $\frac43\varepsilon_0$, which sum to $\frac{1409532425}{1610612736}<\frac9{10}$ (exact rational arithmetic); this gives \eqref{eq:selection-stripflux}. Invariance together with $F_X>0$ gives arrival at $X=\frac54$. For the level, $p_0(\frac54)=\frac{49}{192}$, and
\[
\tfrac12Y^{\top}TY\le\tfrac12E_*=\frac{9R^2}{32\tau}\le\frac9{64KR^3}<\frac1{4K};
\]
this is the point of the energy adaptation: no bound on $\lambda_{\max}(T)$ is needed, because $E_*$ itself is $O(R^2/\tau)$. The remaining adverse level terms are bounded as follows. On the strip $|Y|\le\sqrt{E_*/\tau}=3R/(4\tau)\le3/(8KR^4)$, by \eqref{eq:selection-assume}. So at $X=\frac54$
\[
\Bigl|\tfrac12\bigl(X^2-\tfrac14\bigr)a\cdot Y\Bigr|\le\tfrac{21}{32}R|Y|\le\frac{21}{32K},\qquad
\Bigl|\tfrac12X\,Y^{\top}WY\Bigr|+\Bigl|\tfrac16C[Y,Y,Y]\Bigr|\le\tfrac58R|Y|^2+\tfrac R6|Y|^3\le\frac1{K^2},
\]
and $|rQ(\frac54,Y)|\le\frac12rR\cdot\dist((\frac54,Y),S)^2\le\frac12rR\le3\varepsilon_0$, since $Q(S)=0$, $\nabla Q(S)=0$ and $\|Q\|_{C^2}\le R$. Passing to the exact field changes the level by less than $\eta_R\le\frac1{8192}$. The two fields agree at $S$ by the pins, the $C^1$ distance between their gradient fields on the convex chart $D^+$ is at most $\eta_R$ (Lemma~\ref{lem:selection-robust}), and the segment from $S$ to $(\frac54,Y)$ has length less than $1$. The ledger reserves $\frac3{8192}$ for this term. The remaining terms total $\frac{21}{32K}+\frac1{K^2}+3\varepsilon_0+\frac3{8192}$, and the complete bracket is
\[
\frac{21}{32K}+\frac1{K^2}+3\varepsilon_0+\frac3{8192}+\frac1{4K}=\frac{9869}{16777216}<\frac1{768}
\]
(exact rational arithmetic), which gives \eqref{eq:selection-level}.
\end{proof}

\begin{lemma}[inward capture]\label{lem:selection-capture}
Let $m=2R/\tau$, $\delta_0=1/(512R^4)$, and consider the tapered tube
\[
\mathcal T_{\rm in}=\bigl\{-\tfrac12<X<\tfrac12,\ |Y|\le m\bigl(\tfrac14-X^2\bigr)\bigr\}.
\]
\begin{enumerate}[(i)]
\item The inward half-branch of $S$ enters $\mathcal T_{\rm in}$.
\item (\emph{Lateral flux.}) On $\{|Y|=m(\frac14-X^2)\}$, the function $G=|Y|^2-m^2(\frac14-X^2)^2$ satisfies $\dot G<0$.
\item (\emph{Capture.}) The inward half-branch converges to $M$.
\end{enumerate}
\end{lemma}

\begin{proof}
Two facts are used throughout. First, $Rm=2R^2/\tau\le1/(KR^3)$ and $m\le1$, by \eqref{eq:selection-assume}. Second, since $\nabla Q(M)=\nabla Q(S)=0$ and $\|Q\|_{C^2}\le R$, on the tube
\begin{equation}\label{eq:selection-Qtube}
|\nabla Q(X,Y)|\le R\bigl(\tfrac12-|X|+|Y|\bigr)\le R(2+m)\bigl(\tfrac14-X^2\bigr)\le3R\bigl(\tfrac14-X^2\bigr),
\end{equation}
because $\frac12-|X|=(\frac14-X^2)/(\frac12+|X|)\le2(\frac14-X^2)$ and $|Y|\le m(\frac14-X^2)$. Hence, on the tube, $|F_X-(X^2-\frac14)|\le(\frac14-X^2)(\frac{Rm}2+\frac{Rm^2}8+3rR)$, and
\begin{equation}\label{eq:selection-FXtube}
F_X\le-\bigl(\tfrac14-X^2\bigr)\Bigl(1-\frac{Rm}2-\frac{Rm^2}8-3rR\Bigr)<0 .
\end{equation}

(i) The inward half-branch is the other half of the one-dimensional unstable manifold of $S$, so it leaves $S$ tangent to $-v$, where $v$ is the positive eigenvector of Lemma~\ref{lem:selection-slope}. The proof of that lemma gives $E(v_Y)\le\frac83(R^2/\tau)|v_X|^2$, and $|v_Y|^2\le E(v_Y)/\tau$, so $|v_Y|/|v_X|\le\sqrt{8/3}\,R/\tau<m$. The lateral boundary of the tube leaves $S$ with slope $m$, so the branch starts inside $\mathcal T_{\rm in}$, and by \eqref{eq:selection-FXtube} it moves leftward.

(ii) On $G=0$,
\[
\dot G=2Y\cdot F_Y+4m^2X\bigl(\tfrac14-X^2\bigr)F_X .
\]
The leading term is $-2Y^{\top}TY\le-2\tau m^2(\frac14-X^2)^2$. The adverse terms, using $|Y|=m(\frac14-X^2)$, $|X|\le\frac12$, $\frac14-X^2\le\frac14$ and \eqref{eq:selection-Qtube}, are
\begin{equation}\label{eq:selection-tubeterms}
\begin{aligned}
\bigl|2Y\cdot\tfrac12\bigl(X^2-\tfrac14\bigr)a\bigr|&\le R|Y|\bigl(\tfrac14-X^2\bigr)=mR\bigl(\tfrac14-X^2\bigr)^2,\\
|2Y\cdot XWY|=2|X|\,|Y^{\top}WY|&\le2|X|\,m^2R\bigl(\tfrac14-X^2\bigr)^2\le m^2R\bigl(\tfrac14-X^2\bigr)^2,\\
\bigl|Y\cdot C[Y,Y,\cdot\,]\bigr|&\le m^3R\bigl(\tfrac14-X^2\bigr)^3\le\tfrac14m^3R\bigl(\tfrac14-X^2\bigr)^2,\\
|2Y\cdot r\nabla_YQ|&\le6mrR\bigl(\tfrac14-X^2\bigr)^2,
\end{aligned}
\end{equation}
and dividing the last by the leading term $2\tau m^2(\frac14-X^2)^2$ gives the ratio $3rR/(\tau m)=\frac32r\le\frac32\varepsilon_0<1/(4K)$, using $\tau m=2R$. By the bound on $|F_X|$ above,
\[
\bigl|4m^2X\bigl(\tfrac14-X^2\bigr)F_X\bigr|\le2m^2\bigl(\tfrac14-X^2\bigr)^2\Bigl(1+\frac{Rm}2+\frac{Rm^2}8+3rR\Bigr),
\]
where $3rR\le3\varepsilon_0<\frac1{8192}$. Collecting and dividing by $2\tau m^2(\frac14-X^2)^2$, the adverse ratio is at most
\begin{equation}\label{eq:selection-tubeledger}
\begin{aligned}
&\frac R{2\tau m}+\frac R{2\tau}+\frac{mR}{8\tau}+\frac1{4K}+\frac1\tau\Bigl(1+\frac{Rm}2+\frac{Rm^2}8+\frac1{8192}\Bigr)\\
&\qquad=\frac14+\frac R{2\tau}+\frac{R^2}{4\tau^2}+\frac1{4K}+\frac1\tau\Bigl(\frac{8193}{8192}+\frac{R^2}\tau+\frac{R^3}{2\tau^2}\Bigr),
\end{aligned}
\end{equation}
where $R/(2\tau m)=R/(4R)=\frac14$ and $m=2R/\tau$. Since $\tau\ge2KR^5+2R$ and $K=4096$: $R/(2\tau)\le1/(4K)$, $R^2/(4\tau^2)\le1/(16K^2R^6)$, $1/\tau\le1/(2K)$, and the final bracket is at most $\frac{8193}{8192}+\frac1{2K}+\frac1{8K^2}\le\frac32$. So the total ratio is at most
\[
\frac14+\frac1{4K}+\frac1{16K^2}+\frac1{4K}+\frac1{2K}\Bigl(\frac32+\frac1{2K}+\frac1{8K^2}\Bigr)<\frac14+\frac3K<\frac13 .
\]
Hence $\dot G\le-\frac23\tau m^2(\frac14-X^2)^2<0$ on the whole lateral boundary: trajectories cannot exit through the lateral walls, near $S$ included.

(iii) At $M$ the Jacobian $J_M=\begin{pmatrix}A_M&B_M^{\top}\\B_M&-T+E_D\end{pmatrix}$ is negative definite: for $w=(w_X,w_Y)$,
\[
w^{\top}J_Mw\le-\tfrac34w_X^2+2R|w_X||w_Y|-\tau|w_Y|^2+\tfrac{3R}4|w_Y|^2\le-\tfrac38w_X^2-\tfrac\tau2|w_Y|^2,
\]
by Young's inequality and $\tau\ge2KR^5+2R$. Write $z=(X,Y)$ and $F(z)=J_M(z-M)+\mathcal E(z)$. The polynomial part of $\mathcal E$ is the quadratic part of \eqref{eq:selection-gradient} at $M$, of norm at most $(1+\frac32R)(X+\frac12)^2+2R|Y|^2\le3R|z-M|^2$. Only the $C^2$ norm of $Q$ is controlled, so the $rQ$ part of $\mathcal E$ is bounded linearly: it is $r[\nabla Q(z)-D^2Q(M)(z-M)]$, of norm at most $2rR|z-M|$, because $\nabla Q(M)=0$. Hence the Lyapunov function $V=(X+\frac12)^2+|Y|^2$ satisfies
\[
\dot V=2(z-M)\cdot F(z)\le-\tfrac34V+6RV^{3/2}+4rRV,
\]
which on $\partial B_{\rho_B}(M)$, with $\rho_B:=1/(16R)$, is at most $(-\frac34+\frac38+4\varepsilon_0)V<0$. So the ball is forward invariant, and every trajectory entering it converges to the sink $M$. Take the section $X=-\frac12+2\delta_0$, where $\frac14-X^2=2\delta_0(1-2\delta_0)\le2\delta_0$. There the tube has half-width at most $2m\delta_0=1/(128\,\tau R^3)$, and its points lie within $2\delta_0\sqrt{1+m^2}\le4\delta_0=\frac1{128R^4}<\frac1{16R}=\rho_B$ of $M$. By \eqref{eq:selection-FXtube}, $X$ decreases monotonically along the inward branch inside the tube, so the branch reaches this section inside $B_{\rho_B}(M)$ and is captured.
\end{proof}

\begin{lemma}[exact-field robustness]\label{lem:selection-robust}
Every strict inequality in Lemmas~\ref{lem:selection-slope}--\ref{lem:selection-capture} has a margin of at least $\frac1{10}q^2$, $\frac23q^2$, $\frac13$ of the leading term, or a fixed fraction of $\rho_B^2$, and the adverse brackets are at most $\frac13$, $\frac9{10}$ and $\frac1{768}$ of their budgets. A pin-preserving $C^1$ perturbation of the gradient field of size $\eta_R=1/(8192R^3)$ changes each flux by less than half the margin, so the conclusions of these lemmas persist. For the conditioned side-$24$ field, $\|F_{\rm exact}-F\|_{C^1(D^+)}\le C_+r$ on the fifth-derivative good event, and $C_+r\le\eta_R$ for
\begin{equation}\label{eq:selection-rmat}
r\le r_{\rm mat}=\min\Bigl\{1,\ \bigl[1/(8192\,C_+\varepsilon_0^{3/5})\bigr]^{5/2}\Bigr\}.
\end{equation}
\end{lemma}

\begin{proof}
In $\mathcal C_E$ the perturbation error is $O(\eta_R\xi)$ against $q\ge2R\xi$; on $\{E=E_*\}$ it is at most $q\eta_R\le\frac{4\eta_R}{3R}q^2$; and in the tube it is $O(\eta_R|Y|)$ against the floor $\tau m^2(\frac14-X^2)^2$. In each case this is less than half of the corresponding margin. The bound $\|F_{\rm exact}-F\|_{C^1(D^+)}\le C_+r$ holds on the fifth-derivative good event of Lemma~\ref{lem:selection-fifth}, and \eqref{eq:selection-rmat} gives $C_+r\le\eta_R$.
\end{proof}

\begin{theorem}[energy-adapted capture and escape]\label{thm:selection-escape}
Under \eqref{eq:selection-potential}--\eqref{eq:selection-assume} and $0<r\le r_{\rm mat}$:
\begin{enumerate}[(i)]
\item the $M$-ward unstable half-branch of $S$ converges to $M$ through $\mathcal T_{\rm in}$;
\item the outward half-branch remains in $\mathcal C_E\cup\mathcal S_E$ and reaches $X=\frac54$ at a level strictly above $f(M)$, hence cannot terminate at $M$;
\item both conclusions hold for the exact conditioned field on the fifth-derivative good event.
\end{enumerate}
No bound on $\lambda_{\max}(T)$ is assumed or used.
\end{theorem}

\begin{proof}
Part (i) is Lemma~\ref{lem:selection-capture}; part (ii) is Lemmas~\ref{lem:selection-slope}, \ref{lem:selection-cone} and~\ref{lem:selection-strip}; part (iii) is Lemma~\ref{lem:selection-robust}. Only $\tau=\lambda_{\min}(T)$ enters these proofs.
\end{proof}

Theorem~\ref{thm:selection-escape} is the capture--escape theorem of the main text, with the potential \eqref{eq:selection-potential} as (3.18a) there. The proofs of Lemmas~\ref{lem:selection-cone} and~\ref{lem:selection-capture} contain five local repairs to the proof in the 1~August supplement: the bound on $rQ_X$ in Lemma~\ref{lem:selection-cone}, and, in Lemma~\ref{lem:selection-capture}, the sign of the tangent eigenvector, containment near $S$ through \eqref{eq:selection-Qtube}, the position of the capture section, and the $rQ$ part of the energy remainder. None of them changes a statement.

\subsection{Transverse normalization, the shallow layer and adjacency failures}\label{ssec:selection-adjacency}

\paragraph{Transverse normalization.} The matrix $T$ in \eqref{eq:selection-potential} must be read in scaled coordinates. If $Q_\perp$ is the physical transverse Hessian at the contact saddle, then
\begin{equation}\label{eq:selection-Tr}
T_r=-\frac{Q_\perp}{\kappa r},\qquad\tau_r=\lambda_{\min}(T_r)=\frac{\lambda_*(-Q_\perp)}{\kappa r},
\end{equation}
where $\lambda_*(-Q_\perp)$ is the smallest eigenvalue of $-Q_\perp$. The factor $1/r$ is essential: physical transverse displacement is $rY$, while the fold potential is divided by $\kappa r^3$; hence a physical quadratic term $(r^2/2)\,Y^{\top}Q_\perp Y$ becomes $(1/(2\kappa r))\,Y^{\top}Q_\perp Y$. Equation \eqref{eq:selection-Tr} is (3.18) of the main text.

In the pin-preserving degree-four scaled potential of the conditioned side-$24$ field, the remaining coupling tensors $a,W,C$ and the quartic remainder coefficients $Q_4$ have the polynomial envelope
\begin{equation}\label{eq:selection-RU}
R(U)=1+\|a\|+\|W\|+\|C\|+\|Q_4\|,
\end{equation}
where $U$ collects the random coefficients. By Theorem~\ref{thm:selection-escape}, the conditions
\begin{equation}\label{eq:selection-conditions}
rR(U)^5\le\varepsilon_0,\qquad\tau_r\ge2KR(U)^5+2R(U),
\end{equation}
together with $r\le r_{\rm mat}$ and the fifth-derivative good event, imply simultaneously: (i) the $M$-ward unstable half-branch from $S$ converges to $M$; and (ii) the other half-branch reaches a fixed forward section at field value strictly above $b$ and therefore cannot terminate at $M$. The proof uses $-T_rY$ for transverse contraction, bounds every coupling by $R(U)$, and chooses the cone, strip and tube widths proportional to $R(U)/\tau_r$; condition \eqref{eq:selection-conditions} makes the transverse forcing smaller than the axial margins of $X^2-\frac14$, and the fixed-chart $C^1$ robustness transfers the strict inequalities to the exact field.

\begin{lemma}[shallow layer]\label{lem:selection-shallow}
Because of \eqref{eq:selection-Tr}, failure of the second inequality in \eqref{eq:selection-conditions} inside the typed support is exactly the physical shallow-eigenvalue event
\begin{equation}\label{eq:selection-shallow}
\mathcal S_r=\Bigl\{0<\lambda_*\le\kappa r\bigl[2KR(U)^5+2R(U)\bigr]\Bigr\}\subset\bigl\{0<\lambda_*\le r\,\Psi(U)\bigr\}
\end{equation}
on compact positive $\kappa$ sets, for a fixed polynomial $\Psi$. On the typed support the endpoint determinant product satisfies
\begin{equation}\label{eq:selection-Wshallow}
W_r\le Cr^2\bigl[\lambda_*+r\Psi(U)\bigr]^2(1+\|U\|)^N ,
\end{equation}
and
\begin{gather}
\E^0\bigl[W_r\ind_{\mathcal S_r}\bigr]\le Cr^2\int_0^{r\Psi(U)}\lambda^2\,d\lambda\le Cr^5,\label{eq:selection-shallowmass}\\
\Prob^{MS}_r(\mathcal S_r)\le Cr^3 .\label{eq:selection-shallowP}
\end{gather}
\end{lemma}

\begin{proof}
At contact the transverse Hessian is $-bI_2+{\rm GOE}_2$ (Lemma~\ref{lem:selection-pairlaw}), jointly nondegenerate with the remaining finite jet; side-$24$ jet perturbation and compactness preserve bounded conditional densities and the shallow-eigenvalue coarea estimate. A bounded Gaussian density and integration in the smallest-eigenvalue coordinate, applied to \eqref{eq:selection-Wshallow}, give the unnormalized shallow-layer mass \eqref{eq:selection-shallowmass}. Dividing by $Z_r\ge c_Zr^2$ from \eqref{eq:selection-normalizer} yields \eqref{eq:selection-shallowP}.
\end{proof}

\begin{lemma}[fifth-derivative control]\label{lem:selection-fifth}
Let $H_5$ be the supremum over $\T^3_{24}$ of the norm of the fifth-derivative field of $f$. On compact marks there are constants $C,c>0$ such that
\begin{equation}\label{eq:selection-borell}
\Prob^0(H_5>u)\le C\exp(-cu^2),\qquad u>0 .
\end{equation}
Taking $u=r^{-1}$ gives the failure bound $\exp(-c/r^2)$ for the complement of the \emph{fifth-derivative good event} $\{H_5\le r^{-1}\}$.
\end{lemma}

\begin{proof}
For the fixed side-$24$ periodized Bargmann--Fock field every spectral moment is finite, and every derivative process is continuous on the compact torus. Conditioning on the corrected pair pins changes the fifth-derivative field by a uniformly bounded affine mean and a covariance dominated by the unconditional covariance. On compact marks, the Borell--TIS inequality \cite{AdlerTaylor} therefore gives \eqref{eq:selection-borell}.
\end{proof}

\begin{lemma}[regular-domain and fifth-derivative failures]\label{lem:selection-failure}
Let $F_r$ be the union of the regular-domain failure $\{rR(U)^5>\varepsilon_0\}$ and the fifth-derivative failure $\{H_5>r^{-1}\}$. Then
\begin{equation}\label{eq:selection-failure}
\Prob^{MS}_r(F_r)=\frac{\E^0[W_r\ind_{F_r}]}{Z_r}\le C\,\Prob^0(F_r)^{1/2},
\end{equation}
and $\Prob^{MS}_r(F_r)=o(r^3)$.
\end{lemma}

\begin{proof}
The inequality in \eqref{eq:selection-failure} follows from the exact endpoint second moment $\E^0[W_r^2]\le C_Wr^4$, the Cauchy--Schwarz inequality and $Z_r\ge c_Zr^2$, all from \eqref{eq:selection-normalizer}. The corrected finite-jet regular-domain tail is super-polynomial, and the conditioned fifth-derivative Borell tail is $\exp(-c/r^2)$ by Lemma~\ref{lem:selection-fifth}, so both contributions are $o(r^3)$.
\end{proof}

\begin{proposition}[adjacency and self-attachment]\label{prop:selection-adjacency}
Uniformly for $t\in S^2$ and on compact positive mark sets, for $0<r\le r_{\rm mat}$,
\begin{equation}\label{eq:selection-adjacency}
\Prob^{MS}_r(A_r^c)\le C_Ar^3,\qquad\Prob^{MS}_r(L_{{\rm same},r})\le C_Lr^3 .
\end{equation}
\end{proposition}

\begin{proof}
Outside $\mathcal S_r\cup F_r$ the conditions \eqref{eq:selection-conditions} and the fifth-derivative good event hold, with $T=T_r$. By Theorem~\ref{thm:selection-escape}(i) the $M$-ward half-branch converges to $M$, so $A_r$ holds; by Theorem~\ref{thm:selection-escape}(ii) the outward half-branch cannot terminate at $M$, which excludes same-maximum self-attachment. Hence both probabilities in \eqref{eq:selection-adjacency} are bounded by $\Prob^{MS}_r(\mathcal S_r)+\Prob^{MS}_r(F_r)$, which is $O(r^3)$ by Lemmas~\ref{lem:selection-shallow} and~\ref{lem:selection-failure}.
\end{proof}

\subsection{The selection theorem}\label{ssec:selection-theorem}

\begin{theorem}[elder selection on compact marks]\label{thm:selection-main}
For every compact $I_b\subset\R$ and every compact $I_\kappa\subset(0,\infty)$ there are $r_0>0$ and $C_{\rm sel}<\infty$ such that, under the all-typed maximum--index-two-saddle Palm law \eqref{eq:selection-palm},
\begin{equation}\label{eq:selection-main}
\sup_{t\in S^2,\ b\in I_b,\ \kappa\in I_\kappa}\bigl[1-p_r(t,b,\kappa)\bigr]\le C_{\rm sel}\,r^3,\qquad0<r\le r_0 .
\end{equation}
In particular, the elder-rule selection probability tends to one uniformly on compact positive mark sets.
\end{theorem}

\begin{proof}
By Lemmas~\ref{lem:selection-collar}, \ref{lem:selection-sing} and~\ref{lem:selection-far}, the three regional Palm expectations satisfy
\begin{equation}\label{eq:selection-regional}
\E^{MS}_rN_{\rm col}\le C_{\rm col}r^3,\qquad\E^{MS}_rN_{\rm sing}\le C_{\rm sing}r^3,\qquad\E^{MS}_rN_{\rm far}\le C_{\rm far}r^3,
\end{equation}
so \eqref{eq:selection-markov} gives $\Prob^{MS}_r(B_r)\le C_Br^3$. Proposition~\ref{prop:selection-adjacency} gives $\Prob^{MS}_r(A_r^c)\le C_Ar^3$ and $\Prob^{MS}_r(L_{{\rm same},r})\le C_Lr^3$. By Corollary~\ref{cor:selection-reduction},
\[
1-p_r=\Prob^{MS}_r(E_r^c)\le\Prob^{MS}_r(A_r^c)+\Prob^{MS}_r(E_r^c\cap A_r)\le\Prob^{MS}_r(A_r^c)+\Prob^{MS}_r(B_r)+\Prob^{MS}_r(L_{{\rm same},r})\le C_{\rm sel}r^3,
\]
with all constants uniform in $t\in S^2$ and on the compact mark set.
\end{proof}

Theorem~\ref{thm:selection-main} is the bound (3.13) of the main text, proved here through explicit witnesses in $d=3$.

\begin{remark}[the parent theorem]\label{rem:selection-parent}
The bound \eqref{eq:selection-main} is also Theorem~A of \cite{parentP} (Theorem~\ref{thm:parent-A}) at $d=3$ and $L=24$, and the main text cites that theorem for (3.13). It is proved there for every fixed $d\ge2$ and $L>0$, uniformly over compact birth and gap windows and over all orthonormal frames, with existential constants. Its selection probability, that the global superlevel elder partner of $M$ is $S$, is the $p_r$ of Definition~\ref{def:selection-prob}; on $\Omega_{\rm Mo}$ the two events agree by Lemma~\ref{lem:selection-mark}. Theorem~\ref{thm:selection-main} is a second, three-dimensional proof that does not use \cite{parentP}.
\end{remark}

\begin{remark}[the Morse event under the Palm law]\label{rem:selection-morse}
Lemma~\ref{lem:selection-trichotomy} is applied under $\Prob^{MS}_r$, so it needs $\Prob^{MS}_r(\Omega_{\rm Mo})=1$. Since $\Prob(\Omega_{\rm Mo})=1$, disintegration over the pin values gives this for almost every $(b,\kappa)$, which suffices for every integrated intensity. It holds for every $(b,\kappa)$ as well. Under the pins, $H_M$ and $H_S$ have nondegenerate Gaussian laws, so $M$ and $S$ are nondegenerate critical points almost surely. On each compact set of points at distance at least $1/n$ from $\{M,S\}$, the conditional density of $\nabla f(z)$ is bounded, by derivative-distribution independence (Section~\ref{sec:setting}), and a Bulinskaya-type lemma (see \cite{AdlerTaylor}, Chapter~11) excludes degenerate critical points there; the union over $n$ excludes them everywhere. Equal critical values are excluded in the same way. On pairs of points at distance at least $1/n$ from each other and from $\{M,S\}$, the field $(z,z')\mapsto(\nabla f(z),\nabla f(z'),f(z)-f(z'))$ has seven components in six variables and a bounded conditional density. For a critical point $z\notin\{M,S\}$ with $f(z)\in\{f(M),f(S)\}$, the map $z\mapsto(\nabla f(z),f(z))$ has four components in three variables.
\end{remark}
